% concatenated from nc_polished_compact.tex

\documentclass[reqno,12pt]{amsart}
\usepackage{bbm}
\usepackage[all,pdf]{xy}
\usepackage{epsfig}
\usepackage{amsmath}
\usepackage{amssymb}
\usepackage{amscd}
\usepackage{graphicx}
\usepackage{color}

\usepackage{mathptmx}

\allowdisplaybreaks[4]

\makeatletter
\@namedef{subjclassname@2020}{%
	\textup{2020} Mathematics Subject Classification}
\makeatother
\usepackage[colorlinks=true]{hyperref}
\usepackage{amsfonts}

\usepackage[textwidth=17cm, textheight=24.5cm, centering]{geometry}
\usepackage[nameinlink,noabbrev,nosort]{cleveref}

\raggedbottom
\setlength{\emergencystretch}{2em}

\overfullrule5pt

\newtheorem{thm}{Theorem}[section]
\newtheorem{lemma}[thm]{Lemma}
\newtheorem{ex}[thm]{Example}
\newtheorem{cl}{Claim}
\newtheorem{prop}[thm]{Proposition}
\newtheorem{con}[thm]{Conjecture}
\newtheorem{fact}{Fact}
\newtheorem{ques}[thm]{Question}
\newtheorem{cor}[thm]{Corollary}
\newtheorem{defn}[thm]{Definition}
\newtheorem{rem}[thm]{Remark}
\newtheorem{step}{Step}

\Crefname{thm}{Theorem}{Theorems}
\Crefname{lemma}{Lemma}{Lemmas}
\Crefname{prop}{Proposition}{Propositions}
\Crefname{subsection}{Subsection}{Subsections}

\newcommand{\norm}[1]{\left\Vert #1\right\Vert}
\newcommand{\nnorm}[1]{\lvert\!|\!| #1|\!|\!\rvert}

\renewcommand{\thefootnote}{$\dagger$}

\def \la {\longrightarrow}
\def \ul {\underline}
\def \ol {\overline}
\def \N {\mathbb N}
\def \C {\mathbb C}
\def \Z {\mathbb Z}
\def \R {\mathbb R}
\def \Q {\mathbb Q}
\def \E {\mathbb E}

\def\A {\mathcal A}
\def \M {\mathcal M}
\def\W {\mathcal W}
\def\U {\mathcal U}
\def\V {\mathcal V}
\def\bu {{\bf U}}
\def\I {\mathcal I}
\def\J {\mathcal J}
\def\B {\mathcal B}
\def\P {\mathcal P}

\def \X {\mathcal{X}}
\def \Y {\mathcal{Y}}
\def \ZZ {\mathcal{Z}}
\def \O {\mathcal{O}}
\def \t {\mathcal{T}}

\def \a {\alpha }
\def \b {\beta}
\def \ep {\epsilon}
\def \d {\delta}
\def \D {\Delta}

\def\lk {\left}
\def\re {\right}

\def\FS {{\bf FS}}
\def\FP {{\bf FP}}

\def\AP {{\bf AP}}
\def\GP {{\bf GP}}
\parskip 1.0ex
\numberwithin{equation}{section}
\def\htop{h_{\rm top}}

\begin{document}
\hypersetup{linkcolor=red}
	
\title[]{Pointwise convergence of double ergodic averages along certain non-polynomial sequences}

	\author[]{Rongzhong Xiao}		
	\address[Rongzhong Xiao]{School of Mathematical Sciences, 
		University of Science and Technology of China, 
		Hefei, Anhui, 230026, PR China}
	\email{xiaorz@mail.ustc.edu.cn}
	
	\subjclass[2020]{Primary: 37A30; Secondary: 37A10.}
	\keywords{Double ergodic averages, Ergodic theorems, Measurable flows, Pointwise convergence}

\begin{abstract}
	Fix $c\in (1,2)$. Let $\alpha$ and $\beta$ be two non-zero real numbers. It is shown that for any measure preserving system $(X,\X,\mu,T)$ and any $f,g\in L^{\infty}(\mu)$, the limit
	\begin{equation*}
	\lim_{N\to\infty}\frac{1}{N}\sum_{n=1}^{N}f(T^{\lfloor \alpha n^c \rfloor}x)g(T^{\lfloor \beta n^c \rfloor}x)
	\end{equation*}
	exists for $\mu$-a.e. $x\in X$.
\end{abstract}

\maketitle 
\section{Introduction}
%\subsection{Background on the double ergodic averages}
By {\bf measure preserving system}, we mean a tuple $(X,\X,\mu,T)$, where $(X,\X,\mu)$ is a Lebesgue probability space and $T:X\to X$ is an invertible measure preserving transformation. 
%We say that $(X,\X,\mu,T)$ is {\bf ergodic} if the sub-$\sigma$-algebra $\mathcal{I}(T)$ generated by all $T$-invariant subsets is trivival.

Fix a measure preserving system $(X,\X,\mu,T)$. In the study of ergodic theory, a fundamental problem is to understand the pointwise limit behavior of the following double ergodic averages:
\begin{equation}\label{eq1-0}
	\frac{1}{N}\sum_{n=1}^{N}T^{\lfloor P(n)\rfloor}f\cdot T^{\lfloor Q(n)\rfloor}g,
\end{equation}
where $P,Q\in \R[n]$ with $\deg P=\deg Q$ and $f,g\in L^{\infty}(\mu)$.

In 1990, Bourgain established the following pointwise ergodic theorem.
\begin{thm}\label{thm1-1}
	$($\cite[Main Theorem]{Bourgain90}$)$ Let $(X,\X,\mu,T)$ be a measure preserving system. Fix two distinct non-zero integers $a$ and $b$. Then for any $f,g\in L^{\infty}(\mu)$, we have that 
	\begin{equation}\label{eq1-1}
		\lim_{N\to\infty}\frac{1}{N}\sum_{n=1}^{N}f(T^{an}x)g(T^{bn}x)
	\end{equation}
	exists for $\mu$-a.e. $x\in X$. 
%	When $T$ is ergodic, if $f$ or $g$ is orthogonal to the closed subspace spanned by all eigenfunctions with respect to $T$, then the limit function is zero.
\end{thm}
In 2025, Krause extended Bourgain's theorem to linear polynomials with one real coefficient and one rational coefficient. 
\begin{thm}\label{thm1-2}
	$($\cite[Theorem 1.1]{Krause2025}$)$ Let $(X,\X,\mu,T)$ be a measure preserving system. Let $\alpha\in \R\backslash \Q$ and $\beta\in \Q$. Then for any $f,g\in L^{\infty}(\mu)$, we have that 
	$$\lim_{N\to\infty}\frac{1}{N}\sum_{n=1}^{N}f(T^{\lfloor \alpha n\rfloor}x)g(T^{\lfloor \beta n\rfloor}x)$$ exists for $\mu$-a.e. $x\in X$.
\end{thm}
Following \Cref{thm1-1,thm1-2}, the simplest case of \eqref{eq1-0} may be the following double ergodic averages:
\begin{equation}\label{eq1-7}
		\frac{1}{N}\sum_{n=1}^{N}T^{n^2}f\cdot T^{-n^2}g,
\end{equation} 
where $f,g\in L^{\infty}(\mu)$.

Unfortunately, \eqref{eq1-7} seems to be out of reach at this stage due to the absence of distinct degrees. However, in this direction, Daskalakis established a pointwise ergodic theorem for the $\lfloor n^c\rfloor$-analogue of \eqref{eq1-7}.
\begin{thm}\label{thm2-1}
	$(${\em a special case of} \cite[Theorem 1.7]{Daskalakis2025}$)$ Let $c\in (1,23/22)$. Let $(X,\X,\mu,T)$ be a measure preserving system. Then for any $f,g\in L^{\infty}(\mu)$, we have that 
	$$\lim_{N\to\infty}\frac{1}{N}\sum_{n=1}^{N}f(T^{\lfloor n^c\rfloor}x)g(T^{-\lfloor   n^c\rfloor}x)=\lim_{N\to\infty}\frac{1}{N}\sum_{n=1}^{N}f(T^{n}x)g(T^{-n}x)$$ for $\mu$-a.e. $x\in X$.
\end{thm}
In fact, the function $t^c$ can be replaced by a general $c$-regularly varying function (see \cite[Theorem 1.7]{Daskalakis2025}) in Daskalakis's result. Furthermore, Daskalakis's result shows that the double ergodic averages along $(\lfloor n^c\rfloor,-\lfloor n^c\rfloor)$ share similarities with those in \eqref{eq1-1}. Very recently, for the almost everywhere convergence part of \Cref{thm2-1}, Fornal and Krause \cite[a special case of Proposition 1.5]{FornalKrause2026} extended the range of $c$ to $(1,7/6)$. 
%\subsection{Main results}

Here, we establish a result that combines aspects of \Cref{thm1-2,thm2-1}.
\begin{thm}\label{TB}
Fix $c\in (1,2)$. Let $\alpha$ and $\beta$ be two non-zero real numbers. Let $(X,\X,\mu,T)$ be a measure preserving system. Then for any $f,g\in L^{\infty}(\mu)$, the limit
\begin{equation}\label{eq1-3}
	\lim_{N\to\infty}\frac{1}{N}\sum_{n=1}^{N}f(T^{\lfloor \alpha n^c \rfloor}x)g(T^{\lfloor \beta n^c \rfloor}x)
\end{equation}
exists for $\mu$-a.e. $x\in X$. 
\end{thm}
When $c\in (1,23/22)$ and $|\alpha|=|\beta|$, \Cref{TB} follows from \cite[Theorem 1.7]{Daskalakis2025}; if $c\in (1,7/6)$ and $\alpha,\beta \in \{1,-1\}$, it follows from \cite[Proposition 1.5]{FornalKrause2026}.
%For these comparisons, use $\lfloor-u\rfloor=-\lfloor u\rfloor-1$ for
%$u\notin\Z$ and absorb the extra iterate into the corresponding function.
%For each $a>0$, the set $\{n\in\N:an^c\in\Z\}$ has density
%zero by \Cref{lem2-2}, so these exceptions do not affect convergence.
%\subsection{A brief proof outline of \Cref{TB}} 
%We compare the discrete averages with continuous averages along
%the same orbits. Passing to the suspension flow replaces the
%integer iterates by real times. The continuous averages $C_N$
%converge almost everywhere by \Cref{thm2-2} and
%\Cref{lem2-4}. Thus the problem is to show that discrete
%sampling does not change their limit. Equidistribution controls
%the errors introduced by smoothing the lifted functions and by
%passing back to suspension height zero.
%
%For orbitally smooth $F,G$, we establish the quantitative comparison
%\[
%\|D_N^\phi(F,G)-C_N^\phi(F,G)\|_{L^2(\nu)}
%\lesssim_{F,G,\phi,c,\alpha,\beta}N^{-\delta_c},
%\qquad \delta_c>0,
%\]
%for each $\phi\in C_c^\infty((0,\infty))$.
%Poisson summation identifies $C_N^\phi$ with the zero-frequency
%term. After the change of variables $u=(t/N)^c$, the remaining
%terms have phases $-mNu^{1/c}$ and a common product
%$F(S^{\alpha N^cu})G(S^{\beta N^cu})$.
%
%\Cref{thm3-1} controls smooth sums of these terms over
%$m\asymp M$ for $M\le N^\theta$, where $\theta<1$.
%Its proof uses a short translation, Cauchy--Schwarz, and spectral
%representation to reduce the bilinear integral to scalar
%correlations of its kernel. Choosing
%$H=N^{-(1+\theta)/2}$ gives both $NH\to\infty$ and $MH\to0$.
%The latter preserves curvature for unequal summation indices:
%stationary phase then produces exponential sums estimated by
%derivative bounds and arithmetic counting. For equal indices,
%an inverse change of variables and Poisson summation instead
%produce a phase whose derivatives do not vanish, allowing a
%higher-derivative estimate.
%
%For the remaining frequencies, repeated integration by parts
%gives an upper bound
%$O(N^{\theta-r(\theta-c+1)})$ after $r$ integrations.
%We choose
%\[
%c-1<\theta<1
%\]
%and then fix $r$ sufficiently large. Both frequency ranges
%therefore contribute a power saving. The squared errors are
%summable along geometric sequences; approximation by smooth
%time weights and interpolation give the comparison for all
%$N$. Removing the orbital smoothing and returning from the
%suspension then prove \Cref{TB}.

\medskip
 \noindent\textbf{Structure of the paper.}
Section~2 reduces \Cref{TB} to \Cref{thm3-1}. Section~3 proves
\Cref{thm3-1}.
The appendices record some estimates used in the proof of \Cref{thm3-1}.

\medskip
\noindent\textbf{Notation.}\label{sec1:notation}
Write $\N=\{1,2,\ldots\}$ and
$\Z_{\ge0}=\{0,1,\ldots\}$. For $x\in\R$, $\lfloor x\rfloor$
is the greatest integer not exceeding $x$,
$\lceil x\rceil=-\lfloor-x\rfloor$, and
$\{x\}=x-\lfloor x\rfloor$. We use
\[
 \begin{gathered}
 e(t)=\exp(2\pi it),\qquad
 \|t\|_{\R/\Z}=\min_{m\in\Z}|t-m|\quad(t\in\R),\\
 \widehat a(\xi)=\int_\R a(u)e(-\xi u)\,du\quad(a\in L^1(\R),\ \xi\in\R).
 \end{gathered}
\]
The notation $A\lesssim_{d_1,\ldots,d_k}B$ means
$A\le C_{d_1,\ldots,d_k}B$ for a positive constant depending only on
the indicated data; $A\asymp_{d_1,\ldots,d_k}B$ means both inequalities.
The notations $O_{d_1,\ldots,d_k}(B)$ and $o_{N\to\infty}(1)$ have their
usual meanings.

For an open set $\mathcal O\subset\R$ and
$\mathbb K\in\{\R,\C\}$, $C^r(\mathcal O;\mathbb K)$ denotes
functions with continuous derivatives through order $r$ on
$\mathcal O$. If $I=[a,b]\subset\mathcal O$ is a closed bounded
interval, we put
$\|f\|_{C^r(I)}=\max_{0\le j\le r}\sup_I|f^{(j)}|$;
the derivatives are restrictions from $\mathcal O$.
A function in $C_c^\infty(\mathcal O;\mathbb K)$ has support compactly
contained in $\mathcal O$ and is extended by zero when used on $\R$.
For a pointwise-defined function $a:I\to\C$, its total variation and
variation norm are
\[
 \operatorname{Var}_I(a)=
 \sup_{\substack{s\in\N,\ t_0<\cdots<t_s\\t_0,\ldots,t_s\in I}}
       \sum_{j=1}^s|a(t_j)-a(t_{j-1})|,
 \qquad
 \mathcal V_I(a)=\sup_I|a|+\operatorname{Var}_I(a).
\]
The same definition is used on $\R$; variation on a singleton is zero.
%Sums indexed by $I\cap\Z$ include all integer endpoints and are zero
%when that set is empty. 
%Scalar integrals are Lebesgue integrals,
%Hilbert-space-valued integrals are Bochner integrals, and inner
%products are linear in their first variable. A measurable function
%is $1$-bounded when its essential supremum is at most one.

\medskip
\noindent\textbf{Acknowledgements.}
The author is supported by the National Natural Science Foundation of China (No. 12426201), and by the Postdoctoral Fellowship Program of CPSF under Grant Number GZB20260746.

\medskip
\noindent\textbf{Declaration on the use of AI.}
The mathematical results, arguments, and proofs in this paper 
were developed by the author, who takes full responsibility for the final manuscript.
AI tools were used only for language polishing and to assist with literature searches . 
%All 
%mathematical statements and references were independently checked by the authors .
\section{Reduction to \texorpdfstring{\Cref{thm3-1}}{Theorem~\getrefnumber{thm3-1}}}
\begin{defn}
	Let $d\in \N$. A tuple $(X,\X,\mu, (T^{{\bf t}})_{{\bf t}\in \R^{d}})$ is called a {\bf measurable flow} if $(X,\X,\mu )$ is a Lebesgue probability space, $(T^{\bf t})_{{\bf t}\in \R^{d}}$ is a $d$-parameter group of invertible measure preserving transformations acting on $X$, and the mapping $$\R^d\times X\rightarrow X, \ ({\bf t}, x)\mapsto T^{{\bf t}}x$$ is measurable.
\end{defn}
Fix a measure preserving system $(X,\X,\mu,T)$, $1$-bounded $f,g\in L^{\infty}(\mu)$, $c\in (1,2)$, and two non-zero real numbers $\alpha,\beta$. Set
\[
 X_0=\bigcap_{j\in\Z}
 \{x\in X:|f(T^jx)|\le1\text{ and }|g(T^jx)|\le1\}.
\]
The set $X_0$ is $T$-invariant and co-null. Redefining $f$ and $g$ to be
zero outside $X_0$, we may assume that $|f(x)|,|g(x)|\le1$ for every $x\in X$.

First, let us construct a measurable flow based on $(X,\X,\mu,T)$.  Let $Y=X\times [0,1)$, $\nu=\mu\times \text{Leb}_{[0,1)}$, and $\Y=\X\otimes \mathcal{B}([0,1))$. For any $t\in \R$, we define $$S^{t}:Y\to Y,(x,s)\mapsto (T^{\lfloor t+s\rfloor}x,t+s\mod{1}).$$ This defines a measurable flow ${\bf Y}:=(Y,\Y,\nu,(S^t)_{t\in \R})$. For any $(x,s)\in Y$, let $\tilde{f}(x,s)=f(x)$ and $\tilde{g}(x,s)=g(x)$. 

Next, we introduce the following averages:
$$D_{N}(\tilde{f},\tilde{g})(y)=\frac{1}{N}\sum_{n=1}^{N}\tilde{f}(S^{\alpha n^c}y)\tilde{g}(S^{\beta n^c}y);$$
$$D_{N}(f,g)(x)=\frac{1}{N}\sum_{n=1}^{N}f(T^{\lfloor \alpha n^c \rfloor}x)g(T^{\lfloor \beta n^c \rfloor}x);$$
$$C_N(\tilde{f},\tilde{g})(y)=\frac{1}{N}\int_{0}^{N}\tilde{f}(S^{\alpha t^c}y)\tilde{g}(S^{\beta t^c}y)dt.$$


\subsection{Lift the convergence to flows}
\begin{lemma}\label[lemma]{lem2-1}
	If for $\nu$-a.e. $y\in Y$, the limit $$\lim_{N\to\infty}D_{N}(\tilde{f},\tilde{g})(y)$$ exists, then for $\mu$-a.e. $x\in X$, the limit 
	$$\lim_{N\to\infty}D_{N}(f,g)(x)$$
	exists.
\end{lemma}
The proof of \Cref{lem2-1} uses \Cref{lem2-2}.
\begin{lemma}\label[lemma]{lem2-2}
	$($\em{Cf.} \cite[Theorem 1.3]{Boshernitzan1994}$)$ Let $c\in (1,2)$. Then for any non-zero real number $\alpha$, the sequence $\alpha n^c$ is equidistributed modulo one. 
\end{lemma}
\begin{proof}[Proof of \Cref{lem2-1}]
For each $s\in[0,1)$ and $z\in\R$, the integers $\lfloor z+s\rfloor$
and $\lfloor z\rfloor$ agree unless $\{z\}\in[1-s,1)$. Therefore, for
any $(x,s)\in X_0\times[0,1)$, we have that
\begin{align}
 |D_N(\tilde f,\tilde g)(x,s)-D_N(f,g)(x)|
\lesssim \frac1N\sum_{n=1}^N
 \left(1_{[1-s,1)}(\{\alpha n^c\})
       +1_{[1-s,1)}(\{\beta n^c\})\right).
 \label{eq2-add-height-error}
\end{align}
\Cref{lem2-2} implies
\begin{equation}
 \limsup_{N\to\infty}
 |D_N(\tilde f,\tilde g)(x,s)-D_N(f,g)(x)|\lesssim s.
 \label{eq2-add-height-limsup}
\end{equation}
For every $j\ge2$, we can choose $s_j\in(0,1/j)$ and a co-null
measurable set $X_j\subset X_0$ such that $D_N(\tilde f,\tilde g)(x,s_j)$
converges for every $x\in X_j$. Fix $\displaystyle x\in\bigcap_{j\ge2}X_j$.
For a fixed $j>1$,
\begin{align*}
 |D_N(f,g)(x)-D_M(f,g)(x)|
 &\le |D_N(f,g)(x)-D_N(\tilde f,\tilde g)(x,s_j)|\\
 &\hspace{0.5cm}+|D_N(\tilde f,\tilde g)(x,s_j)
             -D_M(\tilde f,\tilde g)(x,s_j)|\\
 &\hspace{1cm}+|D_M(\tilde f,\tilde g)(x,s_j)-D_M(f,g)(x)|.
\end{align*}
The middle term tends to zero as $N,M\to\infty$. Taking the joint
limit superior in the other terms and applying
\eqref{eq2-add-height-limsup}, we obtain
\[
 \limsup_{N,M\to\infty}|D_N(f,g)(x)-D_M(f,g)(x)|\lesssim s_j.
\]
Letting $j\to\infty$ proves \Cref{lem2-1}.
\end{proof}
\subsection{Convergence of the continuous averages}
\begin{lemma}\label[lemma]{lem2-3}
	For $\nu$-a.e. $y\in Y$, the limit $$\lim_{N\to\infty}C_{N}(\tilde{f},\tilde{g})(y)$$ exists.
\end{lemma}
The proof of \Cref{lem2-3} uses \Cref{thm2-2}
and \Cref{lem2-4}.
\begin{thm}\label{thm2-2}
	$($\cite[Theorem 2.9]{HSX2023}$)$ For $\nu$-a.e. $y\in Y$, the limit $$\lim_{N\to\infty}\frac{1}{N}\int_{0}^{N}\tilde{f}(S^{\alpha t}y)\tilde{g}(S^{\beta t}y)dt$$ exists.
\end{thm}
Here, \cite[Theorem 2.9]{HSX2023} is applied to the flow
$(S^{\alpha t})_{t\in\R}$ with the ratio $\beta/\alpha$; see also
\cite[Theorem 8.30]{BLM2012} for the pair like $t,2t$.
\begin{lemma}\label[lemma]{lem2-4}
$($\emph{a special case of }\cite[Lemma 7.2]{FN2023}$)$ 	Let $\delta>0$. Let $f\in L^{\infty}(m_{\R})$, where $m_{\R}$ is the Lebesgue measure on $\R$. If the limit $\displaystyle \lim\limits_{M\to\infty}\frac{1}{M}\int_{0}^{M}f(t)dt$ exists,
then the limit
$$\lim\limits_{M\to\infty}\frac{1}{M}\int_{0}^{M}f(t^\delta)dt$$ exists and the two limits are equal.
\end{lemma}


\begin{proof}[Proof of \Cref{lem2-3}]
\Cref{thm2-2} and \Cref{lem2-4} imply
\Cref{lem2-3}.
\end{proof}
\subsection{Reduction to orbitally smooth functions}
First, we introduce a special function family.
\begin{defn}
	Given $F\in L^{\infty}(\nu)$, it is called \textbf{orbitally smooth} if it has a bounded representative such that
	$t\mapsto F(S^t y)$ is smooth and each derivative is bounded uniformly in $t,y$. We denote all orbitally smooth $1$-bounded measurable functions with respect to $(Y,\Y,\nu,(S^t)_{t\in \R})$ by $C^{\infty}_{1}({\bf Y})$. 
\end{defn}
\begin{lemma}\label[lemma]{lem2-6}
	If for all $F,G\in C^{\infty}_{1}({\bf Y})$, $$\lim_{N\to\infty}\left|D_{N}(F,G)(y)-C_{N}(F,G)(y)\right|=0$$ for $\nu$-a.e. $y\in Y$, then for $\nu$-a.e. $y\in Y$,
	$$\lim_{N\to\infty}\left|D_{N}(\tilde{f},\tilde{g})(y)-C_{N}(\tilde{f},\tilde{g})(y)\right|=0.$$
\end{lemma}
\begin{proof}
Choose a nonnegative function $\kappa\in C_c^\infty((-1,1))$ with
$\displaystyle\int_\R\kappa(t)\,dt=1$. For $0<\delta<1/4$, set
\[
 \kappa_\delta(t)=\delta^{-1}\kappa(t/\delta),\qquad
 F_\delta(y)=\int_\R\kappa_\delta(t)\tilde f(S^t y)\,dt,
 \quad
 G_\delta(y)=\int_\R\kappa_\delta(t)\tilde g(S^t y)\,dt.
\]
Positivity and normalization of the kernel imply that both functions
are $1$-bounded. To verify orbital smoothness, write
\[
 F_\delta(S^v y)
 =\int_\R\kappa_\delta(t-v)\tilde f(S^t y)\,dt.
\]
For each integer $r\ge0$,
\[
 \sup_{v,y}\left|\frac{d^r}{dv^r}F_\delta(S^v y)\right|
 \le\|\kappa_\delta^{(r)}\|_{L^1(\R)}
 =\delta^{-r}\|\kappa^{(r)}\|_{L^1(\R)}.
\]
The same argument applies to $G_\delta$. Thus
$F_\delta,G_\delta\in C_1^\infty({\bf Y})$.

Let $E_\delta=[0,\delta]\cup[1-\delta,1)$. If
$\delta\le s\le1-\delta$ and $|t|<\delta$, then
$\lfloor s+t\rfloor=0$. Consequently,
$F_\delta(x,s)=\tilde f(x,s)$ and
$G_\delta(x,s)=\tilde g(x,s)$ at those points. Then
\begin{equation}
 |F_\delta-\tilde f|(x,s)\lesssim1_{E_\delta}(s),\qquad
 |G_\delta-\tilde g|(x,s)\lesssim1_{E_\delta}(s).
 \label{eq2-add-smoothing-boundary}
\end{equation}
For each fixed $s$ and each nonzero real $\gamma$, \Cref{lem2-2} gives
\begin{equation}
 \lim_{N\to\infty}\frac1N\sum_{n=1}^N
 1_{E_\delta}(\{s+\gamma n^c\})=2\delta.
 \label{eq2-add-boundary-discrete-count}
\end{equation}
Observe that the periodic function $u\mapsto1_{E_\delta}(\{s+\gamma u\})$
has continuous Ces\`aro limit $2\delta$. Applying \Cref{lem2-4}, with the power equal to $c$, gives
\begin{equation}
 \lim_{N\to\infty}\frac1N\int_0^N
 1_{E_\delta}(\{s+\gamma t^c\})\,dt=2\delta.
 \label{eq2-add-boundary-continuous-count}
\end{equation}
By \eqref{eq2-add-smoothing-boundary} and
\eqref{eq2-add-boundary-discrete-count},
\eqref{eq2-add-boundary-continuous-count} with $\gamma=\alpha,\beta$,
\begin{align}
 \limsup_{N\to\infty}
 |D_N(F_\delta,G_\delta)-D_N(\tilde f,\tilde g)|
 &\lesssim\delta,\label{eq2-add-smoothing-discrete}\\
 \limsup_{N\to\infty}
 |C_N(F_\delta,G_\delta)-C_N(\tilde f,\tilde g)|
 &\lesssim\delta.\label{eq2-add-smoothing-continuous}
\end{align}
Choose a countable sequence $\delta_j\downarrow0$. By the hypothesis
of \Cref{lem2-6}, there is a common co-null set on which
\[
 D_N(F_{\delta_j},G_{\delta_j})
 -C_N(F_{\delta_j},G_{\delta_j})\longrightarrow0
 \quad\text{for every }j.
\]
At a point of this set,
\eqref{eq2-add-smoothing-discrete} and \eqref{eq2-add-smoothing-continuous}
give
\[
 \limsup_{N\to\infty}
 |D_N(\tilde f,\tilde g)-C_N(\tilde f,\tilde g)|\lesssim\delta_j
 \quad\text{for every }j.
\]
Letting $j\to\infty$ proves \Cref{lem2-6}.
\end{proof}
\subsection{Reduction to \texorpdfstring{\Cref{thm3-1}}{Theorem~\getrefnumber{thm3-1}}}
To state \Cref{lem2-5}, for a fixed $\phi \in C^{\infty}_{c}((0,\infty))$ and $F,G\in C^{\infty}_{1}({\bf Y})$, define
$$D_{N}^{\phi}(F,G)(y)=\frac{1}{N}\sum_{n\ge 1}\phi(n/N)F(S^{\alpha n^c}y)G(S^{\beta n^c}y);$$
$$C_N^{\phi}(F,G)(y)=\frac{1}{N}\int_{0}^{\infty}\phi(t/N)F(S^{\alpha t^c}y)G(S^{\beta t^c}y)dt.$$
\begin{lemma}\label[lemma]{lem2-5}
	Suppose that for all $F,G\in C^{\infty}_{1}({\bf Y})$ and all $\phi \in C^{\infty}_{c}((0,\infty))$, there is some $\delta_c>0$ such that for each $N>1$, 
	$$\norm{D_{N}^{\phi}(F,G)-C_N^{\phi}(F,G)}_{L^{2}(\nu)}\lesssim_{F,G,\phi,c,\alpha,\beta} N^{-\delta_c}.$$
	Then for all $F,G\in C^{\infty}_{1}({\bf Y})$, $$\lim_{N\to\infty}\left|D_{N}(F,G)(y)-C_{N}(F,G)(y)\right|=0$$ for $\nu$-a.e. $y\in Y$.
\end{lemma}
\begin{proof}[Proof of \Cref{lem2-5}]
Fix $F,G\in C_1^\infty({\bf Y})$. First fix a real number $q>1$ and put
$N_j=\lfloor q^j\rfloor$, omitting the finitely many initial repetitions.
For each fixed $\phi$, the assumed power saving implies
\[
 \sum_j \norm{D_{N_j}^\phi(F,G)-C_{N_j}^\phi(F,G)}_{L^2(\nu)}^2<\infty.
\]
Therefore,
\begin{equation}
 D_{N_j}^\phi(F,G)(y)-C_{N_j}^\phi(F,G)(y)\longrightarrow0
 \quad\text{for almost every }y.
 \label{eq2-add-geometric-smooth}
\end{equation}
For $0<\tau<1/2$, choose $\phi_\tau\in C_c^\infty((0,1))$ satisfying
$0\le\phi_\tau\le1$ and $\phi_\tau=1$ on $[\tau,1-\tau]$. Then:
\begin{align}
 |D_N(F,G)-D_N^{\phi_\tau}(F,G)|&\lesssim\tau+1/N,
 \label{eq2-add-time-cutoff-discrete}\\
 |C_N(F,G)-C_N^{\phi_\tau}(F,G)|&\lesssim\tau.
 \label{eq2-add-time-cutoff-continuous}
\end{align}
Take a countable sequence $\tau_l\downarrow0$, and intersect the co-null
sets from \eqref{eq2-add-geometric-smooth} for these $\phi_{\tau_l}$.
On their intersection,
\[
 \limsup_{j\to\infty}|D_{N_j}(F,G)-C_{N_j}(F,G)|\lesssim\tau_l
 \quad\text{for every }l.
\]
It follows that $D_{N_j}(F,G)-C_{N_j}(F,G)\to0$ there.
If $N_j\le N<N_{j+1}$, then
\[
 |D_N(F,G)-C_N(F,G)|
 \lesssim |D_{N_j}(F,G)-C_{N_j}(F,G)|+\frac{N-N_j}{N}.
\]
Since $N_{j+1}/N_j\to q$, the limit superior of the last fraction is at
most $q-1$. Now carry out the preceding argument for
$q=1+1/i$, $i\ge1$, and intersect the resulting countably many co-null
sets. On that intersection,
\[
 \limsup_{N\to\infty}|D_N(F,G)-C_N(F,G)|\lesssim1/i
 \quad\text{for every }i.
\]
Letting $i\to\infty$ proves \Cref{lem2-5}.
\end{proof}
\Cref{lem2-7} verifies the hypothesis of \Cref{lem2-5}
under the assumption of \Cref{thm3-1}.
\begin{lemma}\label[lemma]{lem2-7}
	Fix $F,G\in C^{\infty}_{1}({\bf Y})$ and $\phi \in C^{\infty}_{c}((0,\infty))$. Assuming \Cref{thm3-1}, there is some $\delta_c>0$ such that for each $N>1$, 
	$$\norm{D_{N}^{\phi}(F,G)-C_N^{\phi}(F,G)}_{L^{2}(\nu)}\lesssim_{F,G,\phi,c,\alpha,\beta} N^{-\delta_c}.$$
\end{lemma}
\begin{proof}[Proof of \Cref{lem2-7}]
%Fix $F,G\in C_1^\infty({\bf Y})$ and
%$\phi\in C_c^\infty((0,\infty))$. All estimates below are for their
%fixed bounded orbitally smooth representatives. The constants may depend
%on $F,G,\phi,c,\alpha,\beta$ and on a fixed derivative order, but not on
%$N$, the Fourier index, or $y$. The assertion is immediate if
%$\phi\equiv0$, so choose fixed numbers $0<a<b<\infty$ such that
%$\operatorname{supp}\phi$ is contained in $(a,b)$.
The proof is divided into four steps.

\medskip
\noindent\textbf{Step 1. Derivation of \eqref{eq2-add-poisson-error}.}
For $N\ge2$ and $y\in Y$, set
\[
 q_{N,y}(t)=\phi(t/N)F(S^{\alpha t^c}y)G(S^{\beta t^c}y),
 \qquad t>0,
\]
and set $q_{N,y}(t)=0$ for $t\le0$. Then
$q_{N,y}\in C_c^\infty(\R)$, its support is contained in some $[aN,bN]\subset (0,\infty)$,
and extension by zero introduces no boundary singularity.
%\medskip
%\noindent\textbf{Step 1. Derivation of \eqref{eq2-add-orbital-derivatives}.}
%For every integer $j\ge0$, write
%\[
% \partial_S^jF(z)=
% \left.\frac{d^j}{ds^j}F(S^s z)\right|_{s=0},
% \qquad
% M_j(F)=\sup_{z\in Y}|\partial_S^jF(z)|,
%\]
%and define $\partial_S^jG$ and $M_j(G)$ in the same way.
%All these constants are finite by orbital smoothness; $M_0(F),M_0(G)$
%are at most one. The group law gives
%\[
% \frac{d^j}{ds^j}F(S^s y)=(\partial_S^jF)(S^s y),
%\]
%and the analogous formula for $G$.
%
%Introduce the scalar orbit function
%\[
% H_y(s)=F(S^{\alpha s}y)G(S^{\beta s}y),\qquad s\in\R.
%\]
%By the product and chain rules,
%\[
% H_y^{(j)}(s)=\sum_{i=0}^j\binom ji\alpha^i\beta^{j-i}
% (\partial_S^iF)(S^{\alpha s}y)
% (\partial_S^{j-i}G)(S^{\beta s}y).
%\]
%In particular, with
%\[
% B_j=\sum_{i=0}^j\binom ji|\alpha|^i|\beta|^{j-i}
% M_i(F)M_{j-i}(G),
%\]
%we have $\sup_{s,y}|H_y^{(j)}(s)|\le B_j$. This estimate also covers
%negative coefficients.
%
%We next keep track of the powers of $t$ produced by composition with
%$t^c$. For every integer $\ell\ge0$, there are constants
%$b_{\ell,j}(c)$, $0\le j\le\ell$, such that
%\[
% \frac{d^\ell}{dt^\ell}H_y(t^c)
% =\sum_{j=0}^\ell b_{\ell,j}(c)t^{cj-\ell}H_y^{(j)}(t^c),
% \qquad t>0.
%\]
%Here $b_{0,0}(c)=1$, and coefficients with an index outside their stated
%range are set equal to zero. For completeness, differentiating one
%summand gives
%\begin{align*}
% \frac{d}{dt}\bigl(t^{cj-\ell}H_y^{(j)}(t^c)\bigr)
% &=(cj-\ell)t^{cj-\ell-1}H_y^{(j)}(t^c)\\
% &\quad+c\,t^{c(j+1)-\ell-1}H_y^{(j+1)}(t^c).
%\end{align*}
%Thus the formula follows by induction with the recurrence
%\[
% b_{\ell+1,j}(c)
% =(cj-\ell)b_{\ell,j}(c)+c\,b_{\ell,j-1}(c).
%\]
%For example, the first two derivatives are
%\begin{align*}
% \frac{d}{dt}H_y(t^c)&=ct^{c-1}H_y'(t^c),\\
% \frac{d^2}{dt^2}H_y(t^c)
% &=c(c-1)t^{c-2}H_y'(t^c)
%   +c^2t^{2c-2}H_y''(t^c).
%\end{align*}
The Leibniz rule and the support length $O(N)$ of $q_{N,y}(t)$ now give
%\begin{align*}
% q_{N,y}^{(r)}(t)
% &=\sum_{\ell=0}^r\sum_{j=0}^\ell
% \binom r\ell b_{\ell,j}(c)N^{-(r-\ell)}\\
% &\qquad\cdot\phi^{(r-\ell)}(t/N)
% t^{cj-\ell}H_y^{(j)}(t^c),\qquad t>0.
%\end{align*}
%For $aN\le t\le bN$ and $0\le j\le\ell\le r$,
%\begin{align*}
% N^{-(r-\ell)}|t^{cj-\ell}|
% &\le C_{a,b,c,r}N^{-(r-\ell)+cj-\ell}\\
% &=C_{a,b,c,r}N^{cj-r}
% \le C_{a,b,c,r}N^{r(c-1)}.
%\end{align*}
%The last inequality uses $j\le r$, $c>0$, and $N\ge1$. It remains valid
%when $cj-\ell$ is negative, since $t/N$ stays in the fixed interval
%$[a,b]\subset(0,\infty)$. There are only finitely many summands for a
%fixed $r$. The bounds for $\phi^{(r-\ell)}$ and $H_y^{(j)}$ therefore
%imply
%\[
% \sup_{y\in Y}\sup_{t\in\R}|q_{N,y}^{(r)}(t)|
% \le C_rN^{r(c-1)}.
%\]
%Multiplication by the support length $(b-a)N$, with that fixed factor
%absorbed into the constant, proves
\begin{equation}
 \sup_{y\in Y}\norm{q_{N,y}^{(r)}}_{L^1(\R)}
 \lesssim_{r,F,G,\phi,c,\alpha,\beta}N^{1+r(c-1)}.
 \label{eq2-add-orbital-derivatives}
\end{equation}
%This includes $r=0$. For each fixed $r$, only orbital derivatives of
%$F,G$ of order at most $r$ have been used.

Write
$I_{N,m}(y)=N^{-1}\widehat q_{N,y}(m)$ for all $m\in \Z$. Note that
\begin{align*}
 \frac1N\int_\R|q_{N,y}(t)e(-mt)|\,dt
 \le\frac1N\int_{aN}^{bN}|\phi(t/N)|\,dt
 =\int_a^b|\phi(u)|\,du<\infty.
\end{align*}
For $m\ne0$, integration by parts once gives
\begin{align*}
 \int_\R q_{N,y}(t)e(-mt)\,dt
 =-\frac1{2\pi im}\int_\R q_{N,y}(t)
       \frac{d}{dt}e(-mt)\,dt
 =\frac1{2\pi im}\int_\R q_{N,y}'(t)e(-mt)\,dt.
\end{align*}
Repeating this identity $r$ times gives
\[
 I_{N,m}(y)=\frac1{N(2\pi im)^r}
 \int_\R q_{N,y}^{(r)}(t)e(-mt)\,dt,
 \qquad m\ne0.
\]
Choose the constant $C_r$ in
\eqref{eq2-add-orbital-derivatives} so that its right-hand side is
$C_rN^{1+r(c-1)}$, and put $A_r=C_r/(2\pi)^r$. Then
\begin{equation}\label{eq2-0}
\sup_{y\in Y}|I_{N,m}(y)|
\le A_r N^{r(c-1)}|m|^{-r},\qquad m\ne0.
\end{equation}
For any fixed integer $r\ge2$,
\begin{align*}
 \sum_{m\ne0}\sup_{y\in Y}|I_{N,m}(y)|
 \le 2A_rN^{r(c-1)}\sum_{m=1}^\infty m^{-r}
 \le 2A_rN^{r(c-1)}\left(1+\frac1{r-1}\right)<\infty.
\end{align*}
%More explicitly, for every positive integer $L$,
%\begin{align*}
% \sup_{y\in Y}\sum_{|m|>L}|I_{N,m}(y)|
% &\le 2A_rN^{r(c-1)}\sum_{m=L+1}^\infty m^{-r}\\
% &\le\frac{2A_r}{r-1}N^{r(c-1)}L^{1-r}
% \longrightarrow0\qquad(L\to\infty).
%\end{align*}
Thus the series is absolutely and uniformly convergent in $y$ for
each fixed $N$. And
\[
 \sum_{m\ne0}\|I_{N,m}\|_{L^2(\nu)}
 \le\sum_{m\ne0}\sup_{y\in Y}|I_{N,m}(y)|<\infty.
\]
The series therefore also converges absolutely in $L^2(\nu)$.
%Its regroupings below, and the corresponding triangle inequalities
%for infinite tails, are justified by these estimates.

For each $y$, $q_{N,y}$ is smooth and compactly supported,
so Poisson summation gives
\[
 \sum_{n\in\Z}q_{N,y}(n)=\sum_{m\in\Z}\widehat q_{N,y}(m).
\]
The left side is a finite sum. Since $q_{N,y}(n)=0$ for $n\le0$, its
quotient by $N$ is $D_N^\phi(F,G)(y)$. The term $m=0$ on the right,
divided by $N$, is $C_N^\phi(F,G)(y)$. Subtracting it yields
\begin{equation}
 D_N^\phi(F,G)-C_N^\phi(F,G)
 =\sum_{m\in\Z\setminus\{0\}}I_{N,m},\qquad
 I_{N,m}(y)=\frac1N\int_\R q_{N,y}(t)e(-mt)\,dt.
 \label{eq2-add-poisson-error}
\end{equation}
The preceding absolute-convergence estimates show that this is also an
identity in $L^2(\nu)$.

\medskip
\noindent\textbf{Step 2. Derivation of \eqref{eq2-add-low-frequency}.}
%\medskip
%\noindent\textbf{Step 3. Change of variables and smooth frequency blocks.}
Set $\rho=1/c$. In the defining integral of $I_{N,m}$, make the
substitution
\[
 u=(t/N)^c,\qquad t=Nu^\rho,\qquad
 \frac{dt}{N}=\rho u^{\rho-1}\,du,\qquad t^c=N^cu.
\]
It follows that
\begin{align}
 I_{N,m}(y)
 &=\int_\R\eta(u)e(-mNu^\rho)
 F(S^{\alpha N^cu}y)G(S^{\beta N^cu}y)\,du,
 \label{eq2-add-poisson-mode}\\
 \eta(u)&=\rho u^{\rho-1}\phi(u^\rho).
 \notag
\end{align}
Here, $\eta$ is defined by the displayed formula for $u>0$ and extended
by zero to $u\le0$. Then we can say that $\eta\in C_c^\infty((0,\infty))$.
%All integrands containing $u^\rho$ are understood as zero outside this
%positive support. Thus the flow scale in \Cref{thm3-1} is exactly
%$R=N^c$.

Choose a nonincreasing function $v\in C^\infty(\R)$ with $0\le v\le1$,
$v(s)=1$ for $s\le5/4$, and $v(s)=0$ for $s\ge7/4$. Put
\[
 \omega(s)=v(s)-v(2s).
\]
This is a real-valued smooth function. It vanishes for $s\le5/8$ and
for $s\ge7/4$, so its support is a compact subset of $(1/2,2)$.
For $J\ge0$ and a positive integer $m$, 
%\begin{align*}
% \sum_{j=0}^J\omega(m/2^j)
% =v(m/2^J).
%\end{align*}
%because $2m\ge2>7/4$. In particular,
\begin{equation}
 v(m/2^J)=\sum_{j=0}^J\omega(m/2^j).
 \label{eq2-add-dyadic-identity}
\end{equation}
%The functions $v,\omega$ are fixed once and for all, independently of $N$.

Choose the fixed parameters
\begin{equation}
 \theta=\frac c2,\qquad \sigma=\frac{1-\theta}{2}=\frac{2-c}{4},
 \qquad d_c=1+\left\lceil\frac2\sigma\right\rceil,
 \qquad P=2^{d_c}.
 \label{eq2-add-fixed-parameters}
\end{equation}
Since $1<c<2$, we have $0<\theta<1$ and $\sigma>0$. For each $N\ge2$,
put $J=\lfloor\theta\log_2N\rfloor$ and $P_0=2^J$. The inequalities
$J\le\theta\log_2N<J+1$ imply
$
N^\theta/2\le P_0\le N^\theta.
$
Hence, every $M=2^j$, $0\le j\le J$, satisfies
$1\le M\le N^\theta$, as required by \Cref{thm3-1}.

%\medskip
%\noindent\textbf{Step 4. Derivation of \eqref{eq2-add-low-frequency}.}
For $0\le j\le J$, define the two finite block sums
\[
 \mathcal B_{N,j}^{+}(y)=\sum_{\ell\ge1}\omega(\ell/2^j)I_{N,\ell}(y),
 \qquad
 \mathcal B_{N,j}^{-}(y)=\sum_{\ell\ge1}\omega(\ell/2^j)I_{N,-\ell}(y).
\]
By \eqref{eq2-add-poisson-mode}, the positive block is
\[
 \mathcal B_{N,j}^{+}(y)=
 \int_\R K_{N,2^j}(u)
 F(S^{\alpha N^cu}y)G(S^{\beta N^cu}y)\,du,
\]
where $K_{N,2^j}$ is the kernel of \Cref{thm3-1}, with the
fixed functions $\eta,\omega$ just chosen. \Cref{thm3-1},
with $R=N^c$
and $\|F\|_\infty,\|G\|_\infty\le1$, gives
\[
 \|\mathcal B_{N,j}^{+}\|_{L^2(\nu)}\lesssim_{c,\phi} N^{-\sigma/(2P)}.
\]
For the negative
block, use that $\omega$ is real to write
\begin{align*}
 \overline{\mathcal B_{N,j}^{-}(y)}
 &=\int_\R\overline{\eta(u)}
 \left(\sum_{\ell\in\Z}\omega(\ell/2^j)e(-\ell Nu^\rho)\right)
 \cdot\overline F(S^{\alpha N^cu}y)
 \overline G(S^{\beta N^cu}y)\,du.
\end{align*}
Apply \Cref{thm3-1} to $\overline\eta,\omega,\overline F,
\overline G$. It gives
\[
 \|\mathcal B_{N,j}^{-}\|_{L^2(\nu)}\lesssim_{c,\phi} N^{-\sigma/(2P)}.
\]
%Then 
%\[
% \|\mathcal B_{N,j}^{+}\|_{L^2(\nu)},\ 
% \|\mathcal B_{N,j}^{-}\|_{L^2(\nu)}
% \le C_*^{1/2}N^{-\sigma/(2P)}.
%\]

\eqref{eq2-add-dyadic-identity} now implies
\[
 \sum_{m\ne0}v(|m|/P_0)I_{N,m}
 =\sum_{j=0}^J\bigl(\mathcal B_{N,j}^{+}+\mathcal B_{N,j}^{-}\bigr).
\]
Put
$\kappa_c=\sigma/(4P)>0$. The triangle inequality yields
\begin{align*}
 \left\|\sum_{m\ne0}v(|m|/P_0)I_{N,m}\right\|_{L^2(\nu)}
 \le\sum_{j=0}^J\left(
 \|\mathcal B_{N,j}^{+}\|_{L^2(\nu)}
 +\|\mathcal B_{N,j}^{-}\|_{L^2(\nu)}\right)
 \lesssim_{c,\phi}(\log N)N^{-\sigma/(2P)}.
\end{align*}
%Since $J+1\le1+\theta\log_2N\lesssim_c1+\log N$, this proves the
%first inequality below. To verify the second, put
%$\kappa_c=\sigma/(4P)>0$. For $N\ge1$,
%\[
% \log N=\int_1^N\frac{dt}{t}
% \le\int_1^N t^{\kappa_c-1}\,dt
% =\frac{N^{\kappa_c}-1}{\kappa_c}.
%\]
%Thus $1+\log N\le(1+\kappa_c^{-1})N^{\kappa_c}$, and
%$-\sigma/(2P)+\kappa_c=-\sigma/(4P)$. We have proved
Then we have 
\begin{align}
 \left\|\sum_{m\ne0}v(|m|/P_0)I_{N,m}\right\|_{L^2(\nu)}
 \lesssim_{c,\phi} N^{-\sigma/(4P)}.
 \label{eq2-add-low-frequency}
\end{align}
%The dependence on $c$ in the last constant is allowed: $\kappa_c$ is
%fixed and positive before $N$ tends to infinity.

\medskip
\noindent\textbf{Step 3. Derivation of \eqref{eq2-add-high-frequency}.}
The remaining part of \eqref{eq2-add-poisson-error} is
\[
 \sum_{m\ne0}\bigl(1-v(|m|/P_0)\bigr)I_{N,m}.
\]
We have $0\le1-v\le1$, and this factor is zero whenever $|m|\le P_0$.
By the absolute $L^2$-convergence proved in {\bf Step 1} and \eqref{eq2-0}, for any fixed integer
$r\ge2$,
\begin{align*}
 &\left\|\sum_{m\ne0}\bigl(1-v(|m|/P_0)\bigr)I_{N,m}
              \right\|_{L^2(\nu)}
\le\sum_{|m|>P_0}\|I_{N,m}\|_{L^2(\nu)}
\le\sum_{|m|>P_0}\sup_{y\in Y}|I_{N,m}(y)|\\
 &\hspace{6cm}\le A_rN^{r(c-1)}\sum_{|m|>P_0}|m|^{-r}
 \le\frac{2A_r}{r-1}N^{r(c-1)}P_0^{1-r}.
\end{align*}

When $r>1$, $P_0\ge N^\theta/2$ gives
\[
 P_0^{1-r}\le(N^\theta/2)^{1-r}
 =2^{r-1}N^{\theta(1-r)}.
\]
%The factor $2^{r-1}$ is absorbed into a constant depending on the fixed
%$r$. Finally,
%\[
% r(c-1)+\theta(1-r)
% =\theta-r(\theta-c+1).
%\]
Consequently, for any fixed integer
$r\ge2$,
\begin{align}
 \left\|\sum_{m\ne0}\bigl(1-v(|m|/P_0)\bigr)I_{N,m}
              \right\|_{L^2(\nu)}
 \lesssim_{F,G,\phi,c,\alpha,\beta,r} N^{\theta-r(\theta-c+1)}.
 \label{eq2-add-high-frequency}
\end{align}
%Dependence on the fixed data $F,G,\phi,c,\alpha,\beta$ is understood in
%this notation, as throughout the proof.

\medskip
\noindent\textbf{Step 4. Fixing the derivative order and concluding.}
From \eqref{eq2-add-fixed-parameters},
\[
 \theta-c+1=1-\frac c2=\frac{2-c}{2}=2\sigma>0.
\]
Choose
\[
 r=\max\left\{2,
 \left\lceil\frac{\theta+\sigma/(4P)}{2\sigma}\right\rceil\right\}
\]
and fix it. Then $r$ depends only on $c$, and satisfies
\[
 \theta-r(\theta-c+1)=\theta-2r\sigma
 \le-\frac{\sigma}{4P}.
\]
\eqref{eq2-add-poisson-error},
\eqref{eq2-add-low-frequency}, and \eqref{eq2-add-high-frequency}
therefore give
\[
 \|D_N^\phi(F,G)-C_N^\phi(F,G)\|_{L^2(\nu)}
 \lesssim_{F,G,\phi,c,\alpha,\beta}N^{-\sigma/(4P)}.
\]
The required exponent is
\begin{equation}
 \delta_c=\frac{\sigma}{4P}
 =\frac{2-c}{16\,2^{d_c}}>0,
 \qquad d_c=1+\left\lceil\frac8{2-c}\right\rceil.
 \label{eq2-add-decay-exponent}
\end{equation}
%Only finitely many orbital derivatives, up to the fixed order $r$, are
%required. Their bounds, the cutoffs, and $r$ do not vary with $N$.
%This proves \Cref{lem2-7} under the assumption of
%\Cref{thm3-1}.
This proves \Cref{lem2-7}.
\end{proof}
Finally, we can prove \Cref{TB}.
\begin{proof}
The conclusion follows from \Cref{lem2-7,lem2-5,lem2-6,lem2-3,lem2-1}.
\end{proof}
% BEGIN REVISED SECTION 3
%\clearpage
\section{Proof of \texorpdfstring{\Cref{thm3-1}}{Theorem~\getrefnumber{thm3-1}}}\label{sec3:main}
The estimate in \Cref{thm3-1} is the analytic input to
\Cref{lem2-7}. In this section, the notation $\|\cdot\|_2$ means the norm in $L^2(\nu)$ unless
another space is specified.
\begin{thm}\label{thm3-1}
Let $\rho,\theta\in(0,1)$, and fix
$\eta\in C_c^\infty((0,\infty);\C)$ and
$\omega\in C_c^\infty((1/2,2);\C)$. Set
\begin{equation}
 \sigma=\frac{1-\theta}{2},\qquad
 d_\theta=1+\left\lceil\frac2\sigma\right\rceil,
 \qquad P=2^{d_\theta}.
 \label{eq3-add-parameters}
\end{equation}
For an integer $N\ge2$ and a real number $M\in[1,N^\theta]$, define
$K_{N,M}:\R\to\C$ by
\begin{equation}
 K_{N,M}(u)=
 \begin{cases}
 \displaystyle\eta(u)\sum_{m\in\Z}\omega(m/M)e(-mNu^\rho),&u>0,\\[1mm]
 0,&u\le0.
 \end{cases}
 \label{eq3-add-kernel}
\end{equation}
For every measurable flow ${\bf Y}=(Y,\Y,\nu,(S^t)_{t\in\R})$,
every $\alpha,\beta\in\R$, every $R>0$, and all $1$-bounded
$F,G\in L^\infty(\nu)$,
\begin{equation}
 \left\|\int_\R K_{N,M}(u)
 F(S^{\alpha Ru})G(S^{\beta Ru})\,du\right\|_2^2
 \lesssim_{\rho,\theta,\eta,\omega}N^{-\sigma/P}.
 \label{eq3-add-main-bound}
\end{equation}
The implicit constant is independent of $N,M,R,\alpha,\beta,{\bf Y},F,G$.
\end{thm}

\noindent\textbf{Proof strategy for \Cref{thm3-1}.}
\Cref{prop3:scalar}, proved by Cauchy--Schwarz and the
spectral representation \cite[Theorem~4.45]{Folland2016}, reduces
\eqref{eq3-add-main-bound} to an integral of the supremum of the absolute value of the
Fourier transform of $K_{N,M}(u+h)\overline{K_{N,M}(u)}$.
We choose $H=N^{\sigma-1}$: the interval $0<h\le H$ satisfies
$NH\to\infty$ and $MH\to0$.
The second condition preserves the required second-derivative bounds
when the product of the two exponential sums is grouped by
$k=m-\ell$.

For $k\ne0$, \Cref{lem3-add-stationary} gives an exponential
sum over the integer index $j$. The derivative estimates in
\Cref{lem3-add-derivatives} and the counting estimates in
\Cref{lem3-add-resonance} give \Cref{prop3:offdiagonal}.

For $k=0$, inverse coordinates and \Cref{lem3-add-dual-formula}
produce a dual exponential sum. Nonvanishing derivatives of the inverse
phase allow \Cref{lem3-add-derivatives}(iii) to be applied, giving
\Cref{prop3:diagonal}. \Cref{sec3-13} integrates
\eqref{eq3-add-offdiagonal-goal}, \eqref{eq3-add-tiny-integral}, and
\eqref{eq3-add-diagonal-high-derivative} to obtain
\eqref{eq3-add-main-bound}. The auxiliary results are collected in
\Crefrange{app:derivatives}{app:resonance}.
%The data $\rho,\theta,\eta,\omega$ are fixed. 
%In the rest of this section, 
%implicit constants may depend
%only on these data or be absolute, unless a further dependence is indicated. For examples,  the implicit constant in $A\lesssim B$ depends
%on $\rho,\theta,\eta,\omega$ or is absolute; the implicit constant in $C\lesssim_{d_1,\ldots,d_k} D$ depends on $\rho,\theta,\eta,\omega,d_1,\ldots,d_k$ or $d_1,\ldots,d_k$.
% In
%particular, a threshold for $N$ may depend on a fixed derivative order,
%but not on $M,h,\xi$, the flow, or its rescaling.

\subsection{Reduction to scalar correlations}\label{sec3:scalar}

If either $\eta$ or $\omega$ is identically zero,
\eqref{eq3-add-main-bound} is immediate. Otherwise,
choose a closed interval $I_0=[a_0,b_0]\subset(0,\infty)$ with
$\operatorname{supp}\eta\subset(a_0,b_0)$.
The sum over $m$ in \eqref{eq3-add-kernel} is finite, and
$K_{N,M}\in C_c^\infty(\R)$.
For real $h,\xi$, define
\begin{equation}
 \mathcal A_h(\xi)=\int_\R
 K_{N,M}(u+h)\overline{K_{N,M}(u)}e(\xi u)\,du.
 \label{eq3-add-correlation}
\end{equation}
This integral is absolutely convergent. Dominated convergence, locally
in $(h,\xi)$, proves continuity, while Cauchy--Schwarz gives
$$|\mathcal A_h(\xi)|\le\|K_{N,M}\|_{L^2(\R)}^2.$$

\begin{prop}\label[prop]{prop3:scalar}
Under the assumptions of \Cref{thm3-1}, for every real
$H\in(0,1]$,
\begin{equation}
 \left\|\int_\R K_{N,M}(u)
 F(S^{\alpha Ru})G(S^{\beta Ru})\,du\right\|_2^2
 \lesssim_{|I_0|}\frac1H\int_0^H
                \sup_{\xi\in\R}|\mathcal A_h(\xi)|\,dh.
 \label{eq3-add-spectral-reduction}
\end{equation}
\end{prop}
\begin{proof}
Choose measurable representatives of $F,G$ bounded by one everywhere,
and put
\[
 V(u)(y)=K_{N,M}(u)F(S^{\alpha Ru}y)G(S^{\beta Ru}y).
\]
Joint measurability of ${\bf Y}$, together with separability of $L^2(\nu)$, implies that $u\mapsto V(u)$ is strongly measurable. Since $$\|V(u)\|_2\le|K_{N,M}(u)|$$ and $K_{N,M}\in C_c^\infty(\R)$,\[
 V\in L^1(\R;L^2(\nu))\cap L^2(\R;L^2(\nu)).
\]
%Changing representatives does not change its integral, by measure
%preservation at each fixed time and Fubini's theorem.

Define
$$V_H(u)=H^{-1}\int_0^H V(u+s)\,ds$$ for almost every $u\in\R$. Then
$$\int_\R V_H(u)\,du=\frac1H\int_0^H\int_\R V(u+s)\,du\,ds=\int_\R V(u)\,du.$$
The support of $V_H$ is contained in $J_H=[a_0-H,b_0]$, whose length is
$$|J_H|=|I_0|+H\le |I_0|+1.$$ 
Note that 
\begin{align*}
& \left\|\int_\R V\right\|_2=\sup_{\substack{W\in L^2(\nu)\\ \|W\|_2=1}} \left|\left\langle\int_{J_H}V_H(u)\,du,W\right\rangle\right|
\\ & \hspace{3cm}\le\sup_{\substack{W\in L^2(\nu)\\ \|W\|_2=1}}
\int_{J_H}|\langle V_H(u),W\rangle|\,du
\le\int_{J_H}\|V_H(u)\|_2\,du.
\end{align*}
Therefore,
\begin{align*}
 &\left\|\int_\R V\right\|_2^2
 \le \left(\int_{J_H}\|V_H(u)\|_2\,du\right)^2
 \le (|I_0|+1)\int_\R\|V_H(u)\|_2^2\,du\\
 &\hspace{3cm}=\frac{|I_0|+1}{H^2}\int_0^H\int_0^H\int_\R
       \langle V(u+s),V(u+s')\rangle\,du\,ds\,ds'.
\end{align*}
%The first inequality is the triangle inequality for Bochner integrals. More explicitly, for every $W\in L^2(\nu)$ with $\|W\|_2=1$,
%\[ 
% \le\int_{J_H}|\langle V_H(u),W\rangle|\,du
% \le\int_{J_H}\|V_H(u)\|_2\,du,
%\]
%and taking the supremum over $W$ gives the first inequality. The second inequality is scalar Cauchy--Schwarz. The final equality follows by expanding the Bochner integral; Fubini is justified by
%$|\langle V(u+s),V(u+s')\rangle|\le |K_{N,M}(u+s)||K_{N,M}(u+s')|$.

Translating $u$ by $s'$ in the innermost integral gives
\[
 \int_\R\langle V(u+s),V(u+s')\rangle\,du=C(s-s'),
 \qquad
 C(h):=\int_\R\langle V(u+h),V(u)\rangle\,du.
\]
For every integrable scalar function $\Phi$ on $[-H,H]$,
\[
 \int_0^H\int_0^H\Phi(s-s')\,ds\,ds'
 =\int_{-H}^H(H-|h|)\Phi(h)\,dh.
\]
Indeed, set $h=s-s'$ and $t=s'$. The square $[0,H]^2$ becomes
\[
 \{(h,t):-H\le h\le H,\ 0\le t\le H,\ 0\le t+h\le H\}.
\]
For fixed $h$, the allowed $t$-interval is
$[\max(0,-h),\min(H,H-h)]$, whose length is $H-|h|$. Applying the above arguments with $\Phi=C$ and using $H-|h|\le H$ gives
\[
 \left\|\int_\R V\right\|_2^2
 \le\frac{|I_0|+1}{H}\int_{-H}^H|C(h)|\,dh.
\]

For a fixed $h\in\R$, let
$F_h=(F\circ S^{\alpha Rh})\overline F$ and
$G_h=(G\circ S^{\beta Rh})\overline G$.
Changing variables by $z=S^{\alpha Ru}y$ in the inner product gives
\begin{align*}
 \langle V(u+h),V(u)\rangle
 ={}&K_{N,M}(u+h)\overline{K_{N,M}(u)}
 \int_Y F_h(z)G_h(S^{(\beta-\alpha)Ru}z)\,d\nu(z).
\end{align*}
Set $U_uW=W\circ S^{(\beta-\alpha)Ru}$ on $L^2(\nu)$.
The last integral is $\langle U_uG_h,\overline{F_h}\rangle$.

We need to verify strong continuity of the action $U$ before applying the spectral theorem. Joint measurability of ${\bf Y}$ implies that
$t\mapsto\langle U_tW,Z\rangle$ is measurable for every $W,Z\in L^2(\nu)$. For
$W\in L^2(\nu)$ and $a\in C_c^\infty(\R)$, define 
\[
 W_a=\int_\R a(t)U_tW\,dt.
\]
The change of variables $r=t+s$ gives
\[
 U_sW_a=\int_\R a(r-s)U_rW\,dr,
\]
and hence
\[
 \|U_sW_a-W_a\|_2
 \le \|a(\cdot-s)-a\|_{L^{1}(\R)}\|W\|_2\longrightarrow0
 \qquad(s\to0).
\]
It remains to show that the span of the vectors $W_a$ is dense in $L^2(\nu)$. Suppose that
$Z$ is orthogonal to every such vector. Then
\[
 0=\langle W_a,Z\rangle
 =\int_\R a(t)\langle U_tW,Z\rangle\,dt
\]
for every $a\in C_c^\infty(\R)$. Thus
$t\mapsto\langle U_tW,Z\rangle$ vanishes almost everywhere for every $W$.
Choose a countable dense subset $\mathcal D\subset L^2(\nu)$ and a single
$t_0$ outside the union of the corresponding null sets. Then
$\langle U_{t_0}W,Z\rangle=0$ for every $W\in\mathcal D$, hence for every
$W\in L^2(\nu)$. Since $U_{t_0}$ is unitary, $Z=0$. Therefore, the span of the vectors $W_a$ is dense in $L^2(\nu)$. This proves the strong continuity of the action $U$.

By the spectral theorem (see \cite[Theorem~4.45]{Folland2016}), there is a
projection-valued measure $E$ on $\R$ such that
$$U_u=\int_\R e(u\xi)\,dE(\xi).$$ For $\varkappa\in L^1(\R)$, we have
\[
 \int_\R\varkappa(u)U_u\,du
 =\int_\R\left(\int_\R\varkappa(u)e(u\xi)\,du\right)dE(\xi).
\]
The operator on the left is defined strongly on each vector. The
spectral functional calculus bounds its norm by the supremum of the
scalar function in parentheses. Apply this to
$\varkappa(u)=K_{N,M}(u+h)\overline{K_{N,M}(u)}$ to obtain
\[
 \left|\int_\R\langle V(u+h),V(u)\rangle\,du\right|
 \le\sup_{\xi\in\R}|\mathcal A_h(\xi)|.
\]
Finally,
$\mathcal A_{-h}(\xi)=e(\xi h)\overline{\mathcal A_h(-\xi)}$,
so the two half-intervals in $h$ give equal suprema.
%Continuity in $\xi$ allows the supremum to be taken over $\Q$;
%it is therefore measurable in $h$. Its bound by
%$\|K_{N,M}\|_{L^2(\R)}^2$ also gives integrability.
This proves \eqref{eq3-add-spectral-reduction}.
%including the case
%$\alpha=\beta$, for which $U$ is the identity group.
\end{proof}

We now fix the real scales
\begin{equation}
 L=N^\sigma,\qquad H=L/N=N^{\sigma-1},\qquad \tau=Nh.
 \label{eq3-add-shift-scales}
\end{equation}
The substitution $h=\tau/N$ in \Cref{prop3:scalar} gives
\begin{equation}
 MH\le N^{-\sigma},\qquad
 \left\|\int V\right\|_2^2
 \lesssim_{|I_0|}\frac1L\int_0^L
       \sup_{\xi\in\R}|\mathcal A_{\tau/N}(\xi)|\,d\tau.
 \label{eq3-add-dual-time-reduction}
\end{equation}
Here $H\to0$, $L\to\infty$, and $MH\to0$ uniformly in the allowed
$M$. These three facts are used separately below.

Let $\psi:(0,\infty)\to\R$ be $\psi(u)=u^\rho$.
For $h\ge0$, put $w_h(u)=\eta(u+h)\overline{\eta(u)}$ on $\R$,
and for $u>0$ define
\begin{equation}
 \Delta_h\psi(u)=\psi(u+h)-\psi(u),\qquad
 \psi_h(u)=\int_0^1\psi'(u+sh)\,ds.
 \label{eq3-add-divided-difference}
\end{equation}
Thus $\Delta_h\psi=h\psi_h$ and $\psi_0=\psi'$.
On every fixed compact subinterval of $(0,\infty)$, all fixed-order
$u$-derivatives of $\Delta_h\psi$ are $O(h)$ for $0\le h\le1$.
Also $\operatorname{supp}w_h\subset\operatorname{supp}\eta$, and
$\|w_h\|_{C^s(\R)}$ is bounded uniformly in $h$ for each fixed $s$.

For each $k\in\Z$, define a smooth function $a_k:\R\to\C$ by
\[
 a_k(j)=\omega(j/M)\overline{\omega((j-k)/M)}\qquad(j\in\R).
\]
For real $j,\xi$ and $0<h\le H$, let
\begin{equation}
 J_{k,j}(\xi,h)=\int_{I_0}w_h(u)
 e\bigl(\xi u-Nk\psi(u)-Nj\Delta_h\psi(u)\bigr)\,du.
 \label{eq3-add-J}
\end{equation}

By the definitions of $K_{N,M}, \mathcal A_h(\xi)$ and \eqref{eq3-add-J}, 
\begin{equation}
 \begin{gathered}
 \mathcal A_h(\xi)=\sum_{k\in\Z}\sum_{j\in\mathcal J_k}
          a_k(j)J_{k,j}(\xi,h),\\
 \mathcal J_k=\{j\in\Z:M/2<j<2M,\ M/2<j-k<2M\}.
 \end{gathered}
 \label{eq3-add-mode-expansion}
\end{equation}
Only $|k|\le2M$ contribute, and $|\mathcal J_k|\le3M$.
Each $\mathcal J_k$ is empty or consists of consecutive integers.
%When applying \Cref{lem3-add-derivatives}, we use the smallest
%closed interval containing $\mathcal J_k$, provided $\mathcal J_k$
%is nonempty. Empty sets contribute zero, and singletons are covered by
%\Cref{lem3-add-derivatives}. Derivatives in $j$ refer to
%$a_k$ and $J_{k,j}$ as functions of a real variable.
Define
\[
 \mathcal D_h(\xi)=\sum_{j\in\mathcal J_0}a_0(j)J_{0,j}(\xi,h),\qquad
 \mathcal O_h(\xi)=\sum_{\substack{k\in\Z\\0<|k|\le2M}}
                       \sum_{j\in\mathcal J_k}a_k(j)J_{k,j}(\xi,h).
\]
Thus $\mathcal A_h=\mathcal D_h+\mathcal O_h$, where $\mathcal D_h$ is the $k=0$ contribution and $\mathcal O_h$ contains the terms with $k\ne0$. 
%We estimate the off-diagonal contribution here; the diagonal contribution is treated in the next subsection.

\subsection{The off-diagonal contribution}\label{sec3:offdiagonal}
In this subsection, we build the following proposition, which includes two estimates of $\mathcal O_h$.
%We estimate the sums over $j$ before summing over $k$. To specify
%constant dependencies throughout the remainder of \Cref{sec3:main},
%write
%\[
% \mathfrak d=(\rho,\theta,\eta,\omega).
%\]
%The interval $I_0$ of \Cref{sec3:scalar} is fixed using only
%$\operatorname{supp}\eta$. For example, if
%$a_\eta=\min\operatorname{supp}\eta$ and
%$b_\eta=\max\operatorname{supp}\eta$, take
%$I_0=[a_\eta/2,2b_\eta]$. All subsequent auxiliary intervals and
%functions are chosen once from the indicated fixed data. A subscript
%$\mathfrak d$ permits dependence on those choices, but not on
%$N,M,h,\xi,R,\alpha,\beta$, the flow, or the normalized inputs.
%Additional dependencies, such as a derivative order or a positive
%exponent, are displayed separately. The appendices specify their own
%independent data sets; they do not inherit $\mathfrak d$.
\begin{prop}\label[prop]{prop3:offdiagonal}
There exist $\tau_*\in(0,1]$ and an integer $N_{\rm off}\ge2$,
depending only on $\mathfrak d:=(\rho,\theta,\eta,\omega)$, such that, for every integer
$N\ge N_{\rm off}$, real $M\in[1,N^\theta]$, and real $h\in(0,H]$,
\begin{align}
 \sup_{\xi\in\R}|\mathcal O_h(\xi)|
 &\lesssim_{\mathfrak d,\epsilon} N^\epsilon
       \left(Mh+\frac1{Nh}+\sqrt{\frac MN}\right)
       &&(\epsilon>0,\tau_*\le Nh\le L),
       \label{eq3-add-offdiagonal-goal}\\
 \sup_{\xi\in\R}|\mathcal O_h(\xi)|
 &\lesssim_{\mathfrak d}
       \sqrt{\frac MN}\min\{M,(Nh)^{-1}\}+MN^{-3/2}
       &&(0<Nh<\tau_*).
       \label{eq3-add-tiny-offdiagonal}
\end{align}
%Estimate~\eqref{eq3-add-offdiagonal-goal} holds for every fixed
%$\epsilon>0$. The threshold $N_{\rm off}$ does not depend on
%$\epsilon$.
\end{prop}
\begin{proof}
The proof is divided into six steps.

\medskip
\noindent\textbf{Step 1. Fixed intervals and a uniform perturbation bound.}
Set
\[
 r_-:=\tfrac12\psi'(b_0),\qquad
 r_+:=2\psi'(a_0),\qquad J_\psi:=[r_-,r_+].
\]
Because $\psi'$ decreases from $+\infty$ to zero, choose
$0<a_1<a_0<b_0<b_1$ with
$\psi'(a_1)>8r_+$ and $\psi'(b_1)<r_-/8$.
Fix these endpoints using only $\rho,I_0$, and put
\[
 I_1=[a_1,b_1],\qquad I_2=[a_1/2,b_1+1].
\]
Thus $I_1+[0,1]\subset I_2$. On $I_2$, the functions
$\psi',-\psi'',\psi^{(3)}$ have positive minima. In particular,
\[
 \begin{gathered}
 d_0:=\min\{r_-,r_+/2\}>0,\qquad
 m_0:=\min_{I_1}|\psi''|>0,\\
 g_0:=\min\{\psi'(a_1)-r_+,r_- -\psi'(b_1)\}>0.
 \end{gathered}
\]
For each integer $s\ge0$, the identity
\eqref{eq3-add-divided-difference}, differentiated in $u$, gives
\[
 \|\Delta_h\psi\|_{C^s(I_1)}
 \le h\max_{\ell\in\{0,\ldots,s\}}\sup_{u\in I_2}|\psi^{(\ell+1)}(u)|
                  \qquad(0\le h\le 1).
\]
Then for $|k|\ge1,j\in \mathcal O_M:=(-1,2M+1),h\in (0,H]$,
\begin{equation}
 \left\|\frac jk\Delta_h\psi\right\|_{C^s(I_1)}
 \le C_{\mathfrak d,s}\frac{Mh}{|k|}
 \le C_{\mathfrak d,s}N^{-\sigma}\qquad (1\le M\le N^{\theta}).
 \label{eq3-add-phase-perturbation}
\end{equation}
Choose one constant $C_{\mathfrak d,8}^*$ valid simultaneously for
$0\le s\le8$, and let
\[
 \delta_0=\frac14\min\{1,d_0,m_0,g_0\},\qquad
 N_{\rm off}=\max\left\{2,
 \left\lceil(C_{\mathfrak d,8}^*/\delta_0)^{1/\sigma}\right\rceil
 \right\}.
\]
Increasing this integer, if necessary, does not change its permitted
dependencies. 
%In particular, $N_{\rm off}$ depends on $\theta$
%through $\sigma=(1-\theta)/2$. 
%The bound $\delta_0\le1/4$ will also
%be used in controlling the location of the critical points.

For the rest of this proof, fix $N,M,h,\xi$ in the stated ranges.
%The variables $k$ and $j$ remain summation indices. 
%Dependence on the
%fixed $N,M,h,\xi$ is suppressed in $u_k(j)$, $D_k(j)$,
%$\Theta_k(j)$, and $b_k(j)$, but the dependence on $k$ is retained.

\medskip
\noindent\textbf{Step 2. Stationary and nonstationary indices.}
For each nonzero integer $k$ with $|k|\le2M$, set
$r_k=\xi/(Nk)\in\R$. For $j\in\mathcal O_M$, define
$p_{k,j}:(0,\infty)\to\R$ by
\begin{equation}
 p_{k,j}(u)=\operatorname{sgn}(k)
 \left(r_k u-\psi(u)-\frac jk\Delta_h\psi(u)\right).
 \label{eq3-add-normalized-phase}
\end{equation}
This is the phase in \eqref{eq3-add-J} divided by $N|k|$.
If $r_k\notin J_\psi$, then
$\operatorname{dist}(r_k,\psi'(I_0))\ge d_0$.
\eqref{eq3-add-phase-perturbation} therefore gives
$|p_{k,j}'|\ge d_0/2$ on $I_0$.
Every derivative of $p_{k,j}$ of order at least two is bounded on $I_1$
independently of $r_k$, since the term $r_ku$ has no such derivative.
For each fixed integer $1\le B\le 7$, the nonstationary estimate
\eqref{eq3-add-nonstationary} yields
\begin{equation}
 J_{k,j}(\xi,h)=O_{\mathfrak d,B}((N|k|)^{-B})
                   \qquad(r_k\notin J_\psi).
 \label{eq3-add-nonstationary-modes}
\end{equation}
%For $B>7$, the additional derivative bounds follow directly from
%\eqref{eq3-add-phase-perturbation}; no further smallness condition is
%needed. 
Taking $B=2$, using $|a_k(j)|\le\|\omega\|_\infty^2$ and
$|\mathcal J_k|\le3M$, gives
\begin{equation}
 \sum_{\substack{k\in\Z,\ 0<|k|\le2M\\r_k\notin J_\psi}}
 \sum_{j\in\mathcal J_k}|a_k(j)J_{k,j}(\xi,h)|
 \lesssim_{\mathfrak d}MN^{-2}.
 \label{eq3-add-nonstationary-total}
\end{equation}
%Here the sum over $k$ is bounded by
%$2\sum_{p\ge1}p^{-2}$, an absolute constant.

Suppose $r_k\in J_\psi$. The function
\[
 \Phi_{k,j}(u)=\psi'(u)+(j/k)(\Delta_h\psi)'(u),\qquad (u>0,j\in \mathcal O_M)
\]
is strictly decreasing on $I_1$: its derivative is at most $-3m_0/4$.
Its value at $a_1$ is at least $r_++3g_0/4$, and its value at $b_1$
is at most $r_- -3g_0/4$. Thus there is a unique
$u_k(j)\in(a_1,b_1)$ satisfying
\begin{equation}
 k\psi'(u_k(j))+j(\Delta_h\psi)'(u_k(j))=\xi/N.
 \label{eq3-add-critical-equation}
\end{equation}
On $I_1$, the derivative upper bound
$|\Phi_{k,j}'|\le B_1:=1+\max_{I_1}|\psi''|$ follows from
$\delta_0\le1/4$. The mean value theorem implies
\[
 u_k(j)-a_1\ge\frac{3g_0}{4B_1},\qquad
 b_1-u_k(j)\ge\frac{3g_0}{4B_1}.
\]
Consequently all these points lie in
\[
 I_*=[a_1+\delta_*,b_1-\delta_*],\qquad
 \delta_*:=\min\{(b_1-a_1)/4,g_0/(4B_1)\}>0.
\]
Note that for any $j\in \mathcal O_M$, we have $$\psi'(u_k(j))+(j/k)(\Delta_h\psi)'(u_k(j))-r_k=0.$$ Therefore, the implicit function theorem gives a smooth function
$u_k:\mathcal O_M\to(a_1,b_1)$; all its values lie in the fixed
interval $I_*$. Define, for real $j\in\mathcal O_M$,
\begin{equation}
 \begin{gathered}
 D_k(j)=k\psi''(u_k(j))+j(\Delta_h\psi)''(u_k(j))=k\Phi_{k,j}'(u_k(j)),\\
 |D_k(j)|\asymp_{\mathfrak d}|k|,\qquad
 \operatorname{sgn}D_k(j)=-\operatorname{sgn}(k).
 \end{gathered}
 \label{eq3-add-critical-curvature}
\end{equation}
The above estimate follows from that on $I_1$, $\Phi_{k,j}'(u_k(j))\le -3m_0/4<0$ and
$$3m_0|k|/4\le |D_k(j)|=|k||\Phi_{k,j}'|\le B_1|k|.$$ 
%The same second-derivative lower bound holds throughout $I_1$, not
%only at $u_k(j)$, because
%$\psi''+(j/k)(\Delta_h\psi)''\le-3m_0/4$ there.

Let
\[
 \mathcal K(\xi)=\{k\in\Z:0<|k|\le2M,\ r_k\in J_\psi\}.
\]
If $\mathcal K(\xi)$ is empty,
\eqref{eq3-add-nonstationary-total} already proves both required
estimates: the quantity $MN^{-2}$ is bounded by both $MN^{-3/2}$
and $\sqrt{M/N}$. Otherwise, define $Q_\xi=|\xi|/N>0$. For $k\in\mathcal K(\xi)$,
\begin{equation}
 \operatorname{sgn}k=\operatorname{sgn}\xi,\qquad
 \frac{Q_\xi}{r_+}\le|k|\le\frac{Q_\xi}{r_-},\qquad
 r_-\le Q_\xi\le2r_+M.
 \label{eq3-add-stationary-k-range}
\end{equation}
%The lower bound uses $|k|\ge1$. 
Hence,
\begin{equation}
 |\mathcal K(\xi)|\lesssim_{\mathfrak d}Q_\xi,\qquad
 \sum_{k\in\mathcal K(\xi)}|k|^{-1}\lesssim_{\mathfrak d}1,
 \qquad |k|\asymp_{\mathfrak d}Q_\xi.
 \label{eq3-add-stationary-count}
\end{equation}
Indeed, the interval for $|k|$ has length
$Q_\xi(r_-^{-1}-r_+^{-1})$, so its number of integers is at most
that length plus one. Each reciprocal is at most $r_+/Q_\xi$;
using $Q_\xi\ge r_-$ gives the second estimate. The sign condition
in \eqref{eq3-add-stationary-k-range} makes $k\mapsto|k|$ injective.

\medskip
\noindent\textbf{Step 3. Stationary phase and the sum over $j$.}
For $k\in\mathcal K(\xi)$ and $j\in\mathcal O_M$, set
\begin{equation}
 \Theta_k(j)=\xi u_k(j)-Nk\psi(u_k(j))
                    -Nj\Delta_h\psi(u_k(j)).
 \label{eq3-add-action}
\end{equation}
The normalized phases $p_{k,j}$ have uniformly bounded $C^8(I_1)$-norms because $r_k\in J_\psi$. Their second derivatives have sign
$\operatorname{sgn}(k)$ and absolute value at least $3m_0/4$ on
$I_1$. The amplitude $w_h$ has support in $I_0$
and uniformly bounded $C^6(\R)$-norm. \Cref{lem3-add-stationary}
therefore applies with phase parameter $N|k|$ and gives
\begin{equation}
 \begin{split}
 a_k(j)J_{k,j}(\xi,h)
 &=(N|k|)^{-1/2}b_k(j)e(\Theta_k(j))+\mathcal R_{k,j},\\
 |\mathcal R_{k,j}|&\le C_{\mathfrak d}(N|k|)^{-3/2}
                         \qquad(j\in\mathcal J_k),
 \end{split}
 \label{eq3-add-mode-stationary}
\end{equation}
where
\begin{equation}
 b_k(j)=e(\operatorname{sgn}(k)/8)
 \frac{a_k(j)w_h(u_k(j))}
 {\bigl|\psi''(u_k(j))+(j/k)(\Delta_h\psi)''(u_k(j))\bigr|^{1/2}}.
 \label{eq3-add-stationary-amplitude}
\end{equation}
The function $b_k$ is smooth on $\mathcal O_M$. 
%The factor
%$e(\operatorname{sgn}(k)/8)$ follows from the sign of $p_{k,j}''$.
%\Cref{lem3-add-stationary} also covers critical points outside
%$\operatorname{supp}w_h$, where its leading coefficient is zero.

Differentiating \eqref{eq3-add-critical-equation} in the real variable
$j$, with $N,M,h,\xi,k$ fixed, gives
\begin{equation}
 u_k'(j)=-\frac{(\Delta_h\psi)'(u_k(j))}{D_k(j)}
                =O_{\mathfrak d}(h/|k|).
 \label{eq3-add-critical-point-derivative}
\end{equation}
To differentiate the denominator in
\eqref{eq3-add-stationary-amplitude}, set $E_k(j)=D_k(j)/k<0$.
The chain rule gives
\[
 E_k'(j)=\left(\psi^{(3)}(u_k(j))
       +\frac jk(\Delta_h\psi)^{(3)}(u_k(j))\right)u_k'(j)
       +\frac1k(\Delta_h\psi)''(u_k(j)).
\]
The parenthesis is bounded by some $C_{\mathfrak d}$; thus
$|E_k'|\le C_{\mathfrak d}h/|k|$. The other factors satisfy
\begin{align*}
 a_k'(j)&=M^{-1}\omega'(j/M)\overline{\omega((j-k)/M)}
   +M^{-1}\omega(j/M)\overline{\omega'((j-k)/M)},\\
 \frac d{dj}w_h(u_k(j))&=w_h'(u_k(j))u_k'(j),\\
 \frac d{dj}(-E_k(j))^{-1/2}
   &=\frac12(-E_k(j))^{-3/2}E_k'(j).
\end{align*}
In particular, $|a_k'|\le2\|\omega\|_\infty\|\omega'\|_\infty/M$.
Using $-E_k\ge3m_0/4$ in the product rule proves
\begin{equation}
 \begin{gathered}
 \|b_k\|_{L^\infty([0,2M])}\le C_{\mathfrak d},\qquad
 |b_k'(j)|\le C_{\mathfrak d}(M^{-1}+h/|k|),\\
 \mathcal V_{[0,2M]}(b_k)
 \le C_{\mathfrak d}\left(1+\int_0^{2M}(M^{-1}+h/|k|)\,dj\right)
 \lesssim_{\mathfrak d}1+Mh/|k|\lesssim_{\mathfrak d}1.
 \end{gathered}
 \label{eq3-add-amplitude-variation}
\end{equation}

%We keep the stationary remainders separate from the leading sums.
Define
\[
 \begin{aligned}
 \mathcal S_k(\xi,h)&=(N|k|)^{-1/2}
             \sum_{j\in\mathcal J_k}b_k(j)e(\Theta_k(j)),\\
 \mathcal E_{\rm ns}(\xi,h)&=
 \sum_{\substack{k\in\Z,\ 0<|k|\le2M\\r_k\notin J_\psi}}
 \sum_{j\in\mathcal J_k}a_k(j)J_{k,j}(\xi,h).
 \end{aligned}
\]
An empty $\mathcal J_k$ contributes zero. Let
$$\mathcal E_{\rm st}(\xi,h)=\sum_{k\in\mathcal K(\xi)}\sum_{j\in\mathcal J_k}\mathcal R_{k,j}.$$ From
\eqref{eq3-add-mode-stationary} and
\eqref{eq3-add-nonstationary-total},
\begin{equation}
 \begin{split}
 \mathcal O_h(\xi)&=\sum_{k\in\mathcal K(\xi)}\mathcal S_k(\xi,h)
                 +\mathcal E_{\rm st}(\xi,h)+\mathcal E_{\rm ns}(\xi,h),\\
 |\mathcal E_{\rm st}(\xi,h)|
 &\le\sum_{k\in\mathcal K(\xi)}\sum_{j\in\mathcal J_k}
                  |\mathcal R_{k,j}|\lesssim_{\mathfrak d}MN^{-3/2},\\
 |\mathcal E_{\rm ns}(\xi,h)|&\lesssim_{\mathfrak d}MN^{-2}.
 \end{split}
 \label{eq3-add-remainder-sum}
\end{equation}
For the stationary bound we used
$|\mathcal J_k|\le3M$ and
$$\sum_{k\ne0}|k|^{-3/2}=2\sum_{p\ge1}p^{-3/2}<\infty.$$
%No cancellation between remainders is needed.

Differentiation of \eqref{eq3-add-action} yields
\[
 \Theta_k'(j)=-N\Delta_h\psi(u_k(j))
 +\bigl(\xi-Nk\psi'(u_k(j))-Nj(\Delta_h\psi)'(u_k(j))\bigr)u_k'(j).
\]
The parenthesis is zero by \eqref{eq3-add-critical-equation}.
Differentiating again and using
\eqref{eq3-add-critical-point-derivative}, we obtain
\begin{align}
 \Theta_k'(j)&=-N\Delta_h\psi(u_k(j)),
       \label{eq3-add-action-first}\\
 \Theta_k''(j)&=N\frac{((\Delta_h\psi)'(u_k(j)))^2}{D_k(j)}.
       \label{eq3-add-action-second}
\end{align}
Since $u_k(j)+[0,h]\subset I_2$ and $\psi''$ has one negative sign,
\[
 h\min_{I_2}|\psi''|
 \le\left|\int_0^h\psi''(u_k(j)+s)\,ds\right|
 =|(\Delta_h\psi)'(u_k(j))|
 \le h\max_{I_2}|\psi''|.
\]
Define the positive scale $\Lambda_\xi=Nh^2/Q_\xi$. Then
\begin{equation}
 |\Theta_k''(j)|\asymp_{\mathfrak d}\Lambda_\xi,
 \qquad \operatorname{sgn}\Theta_k''(j)=-\operatorname{sgn}k
                         \quad(0\le j\le2M).
 \label{eq3-add-mode-curvature}
\end{equation}
%The phase is smooth on $\mathcal O_M\supset[0,2M]$.
\Cref{lem3-add-derivatives}(ii), \eqref{eq3-add-amplitude-variation}, and Abel summation
therefore give
\[
 \left|\sum_{j\in\mathcal J_k}b_k(j)e(\Theta_k(j))\right|
 \lesssim_{\mathfrak d}M\Lambda_\xi^{1/2}+\Lambda_\xi^{-1/2}.
\]
%For a nonempty $\mathcal J_k$, the summation interval is its smallest
%closed containing interval, of length at most $2M$; singleton and empty
%sets obey the same bound. 
Multiplication by $(N|k|)^{-1/2}$ and
$|k|\asymp_{\mathfrak d}Q_\xi$ give
\begin{equation}
 \begin{split}
 |\mathcal S_k(\xi,h)|&\lesssim_{\mathfrak d}
                    \frac{Mh}{|k|}+\frac1{Nh},\\
 \left|\sum_{j\in\mathcal J_k}a_k(j)J_{k,j}(\xi,h)\right|
 &\lesssim_{\mathfrak d}
       \frac{Mh}{|k|}+\frac1{Nh}+M(N|k|)^{-3/2}.
 \end{split}
 \label{eq3-add-mode-second-bound}
\end{equation}
%These inequalities hold before any restriction on the size of
%$\Lambda_\xi$ or the first-derivative error.

For the error $\Theta_k'(j)+{h\xi}/{k}$, the exact identities
\begin{align}
 h\psi'(u)-\Delta_h\psi(u)
     &=-\int_0^h(h-s)\psi''(u+s)\,ds,
          \label{eq3-add-divided-identities}\\
 (\Delta_h\psi)'(u)&=\int_0^h\psi''(u+s)\,ds \nonumber
\end{align}
imply bounds $C_{\mathfrak d}h^2$ and $C_{\mathfrak d}h$, respectively,
on $I_1$. Multiplying \eqref{eq3-add-critical-equation} by $Nh/k$
and using \eqref{eq3-add-action-first}, we obtain
\[
 \Theta_k'(j)+\frac{h\xi}{k}
 =N\left[h\psi'(u_k(j))-\Delta_h\psi(u_k(j))
       +\frac jk h(\Delta_h\psi)'(u_k(j))\right].
\]
Since $|j/k|\le C_{\mathfrak d}M/Q_\xi$ and
$Q_\xi\le2r_+M$, choose $C_U=C_U(\mathfrak d)\ge1$ such that
\begin{equation}
 U_\xi=C_U M\Lambda_\xi=C_U\frac{NMh^2}{Q_\xi},\qquad
 \left|\Theta_k'(j)+\frac{h\xi}{k}\right|\le U_\xi
                      \quad(0\le j\le2M).
 \label{eq3-add-derivative-window}
\end{equation}
%Although $Q_\xi,\Lambda_\xi,U_\xi$ also depend on the fixed
%$N,M,h$, they do not depend on $k$ or $j$. This is why one common
%case distinction applies to the sum over $\mathcal K(\xi)$.

Choose fixed $0<a_\psi\le b_\psi$ such that
$a_\psi\le\psi_h(u)\le b_\psi$ on $I_1$ for $0\le h\le1$.
%One may take the minimum and maximum of $\psi'$ on $I_2$.
Set
\[
 \tau_*:=\min\{1,(4b_\psi)^{-1}\},\qquad
 u_0:=\min\{2^{-12},C_U2^{-20}\}.
\]
These constants depend only on $\mathfrak d$, and
$8u_0+8\sqrt{u_0/C_U}<1/8$.

\medskip
\noindent\textbf{Step 4. The range $\tau=Nh\ge\tau_*$ with $U_\xi\ge u_0$.}
Using \eqref{eq3-add-stationary-count} in the first line of
\eqref{eq3-add-mode-second-bound}, we get
\[
 \sum_{k\in\mathcal K(\xi)}|\mathcal S_k(\xi,h)|
 \lesssim_{\mathfrak d}Mh+\frac{Q_\xi}{Nh}.
\]
The assumption $U_\xi\ge u_0$ implies
$Q_\xi\le(C_U/u_0)NMh^2$, so the second term is at most
$(C_U/u_0)Mh$. Moreover
\[
 MN^{-3/2}\le\sqrt{M/N},\qquad MN^{-2}\le\sqrt{M/N}
                                  \quad(1\le M\le N^{\theta}).
\]
Adding \eqref{eq3-add-remainder-sum} proves
\eqref{eq3-add-offdiagonal-goal} in this case, with no $N^\epsilon$
loss needed.

\medskip
\noindent\textbf{Step 5. The range $\tau=Nh\ge\tau_*$ with $0<U_\xi<u_0$.}
Put
\[
 A_\xi=h|\xi|=\tau Q_\xi,\qquad
 \vartheta_\xi=8U_\xi+8\sqrt{\Lambda_\xi}.
\]
Then $U_\xi\le\vartheta_\xi/8$ and
$0<\vartheta_\xi<1/8$, since
$\Lambda_\xi=U_\xi/(C_UM)$ and $M\ge1$.
For positive integers $p$, define
$d_\xi(p)=\|A_\xi/p\|_{\R/\Z}$.
We apply \Cref{lem3-add-resonance} on
$[Q_\xi/r_+,Q_\xi/r_-]\cap\N$, with
$c_1=1/r_+$, $c_2=1/r_-$, $q_*=r_-$, and $B=1$.
Since $Q_\xi\ge r_-$, $\tau\ge\tau_*$, and
\begin{equation}
 A_\xi+Q_\xi\le2r_+(L+1)M
       \le4r_+N^{\sigma+\theta}\le C_*N,
 \qquad C_*:=\max\{1,4r_+\},
 \label{eq3-add-resonance-size}
\end{equation}
%All the fixed inputs to \Cref{lem3-add-resonance} are thus
%functions of $\mathfrak d$.
the assumptions of \Cref{lem3-add-resonance} are satisfied.

Let
\[
 \mathcal K_{\rm res}(\xi,h)
   =\{k\in\mathcal K(\xi):d_\xi(|k|)\le2\vartheta_\xi\},
 \qquad
 \mathcal K_{\rm nr}(\xi,h)=\mathcal K(\xi)\setminus\mathcal K_{\rm res}(\xi,h).
\]
Using the threshold $2\vartheta_\xi<1/4$ and the injectivity of
$k\mapsto|k|$, \Cref{lem3-add-resonance} gives, for every
$\epsilon'>0$,
\begin{equation}
 |\mathcal K_{\rm res}(\xi,h)|
 \lesssim_{\mathfrak d,\epsilon'}N^{\epsilon'}
                                      (1+Q_\xi\vartheta_\xi).
 \label{eq3-add-resonant-count-used}
\end{equation}
The first line of \eqref{eq3-add-mode-second-bound} now yields
\begin{equation}
 \begin{split}
 \left|\sum_{k\in\mathcal K_{\rm res}(\xi,h)}\mathcal S_k(\xi,h)\right|
 &\lesssim_{\mathfrak d,\epsilon'}N^{\epsilon'}
        (1+Q_\xi\vartheta_\xi)
        \left(\frac{Mh}{Q_\xi}+\frac1{Nh}\right)\\
 &\lesssim_{\mathfrak d,\epsilon'}N^{\epsilon'}
        \left(Mh+\frac1{Nh}+\sqrt{Q_\xi/N}\right).
 \end{split}
 \label{eq3-add-resonant-total}
\end{equation}
To obtain the second line, expand the product into
\[
 \frac{Mh}{Q_\xi}+\frac1{Nh}+Mh\vartheta_\xi
                            +\frac{Q_\xi\vartheta_\xi}{Nh}.
\]
The first and third terms are bounded by $C_{\mathfrak d}Mh$.
For the fourth, use the two exact computations
\[
 \frac{Q_\xi U_\xi}{Nh}=C_U Mh,
 \qquad
 \frac{Q_\xi\sqrt{\Lambda_\xi}}{Nh}=\sqrt{Q_\xi/N}.
\]

For $k\in\mathcal K_{\rm nr}(\xi,h)$,
$d_\xi(|k|)>2\vartheta_\xi$. The distance from $\Z$ is a
$1$-Lipschitz function, and $h\xi/k=A_\xi/|k|$, so
\[
 \operatorname{dist}(\Theta_k'(j),\Z)
 \ge d_\xi(|k|)-U_\xi\ge d_\xi(|k|)/2
                              \quad(0\le j\le2M).
\]
The derivative $\Theta_k'$ is continuous and monotone by
\eqref{eq3-add-mode-curvature}. Its image therefore lies in a
single component of $\R\setminus\Z$. With $\ell_k=\lfloor\Theta_k'(0)\rfloor\in\Z$, we have
\[
 \Theta_k'([0,2M])\subset
 [\ell_k+d_\xi(|k|)/2,\,\ell_k+1-d_\xi(|k|)/2].
\]
%The integer $\ell_k$ is distinct from the real ratio $r_k$ in
%\eqref{eq3-add-normalized-phase}. 
Applying
\Cref{lem3-add-derivatives}(i),
\eqref{eq3-add-amplitude-variation}, and Abel summation gives
\[
 \left|\sum_{j\in\mathcal J_k}b_k(j)e(\Theta_k(j))\right|
       \le C_{\mathfrak d}d_\xi(|k|)^{-1}.
\]
Use the reciprocal-distance estimate of
\Cref{lem3-add-resonance} and
$(N|k|)^{-1/2}\asymp_{\mathfrak d}(NQ_\xi)^{-1/2}$ to conclude
\begin{equation}
 \left|\sum_{k\in\mathcal K_{\rm nr}(\xi,h)}\mathcal S_k(\xi,h)\right|
 \lesssim_{\mathfrak d,\epsilon'}
 \frac{N^{\epsilon'}}{\sqrt{NQ_\xi}}
       \left(\vartheta_\xi^{-1}
                     +Q_\xi\log(2/\vartheta_\xi)\right).
 \label{eq3-add-nonresonant-before}
\end{equation}
Since $\vartheta_\xi\ge8\sqrt{\Lambda_\xi}$,
\[
 \frac1{\sqrt{NQ_\xi}\,\vartheta_\xi}\le\frac1{8Nh}.
\]
For the logarithm term, \eqref{eq3-add-stationary-k-range} and
$\tau\ge\tau_*$ give
\[
 \sqrt{\Lambda_\xi}=\frac{\tau}{\sqrt{NQ_\xi}}
 \ge\frac{\tau_*}{\sqrt{2r_+}}N^{-(1+\theta)/2}
 \ge\frac{\tau_*}{\sqrt{2r_+}}N^{-1}.
\]
Thus $\log(2/\vartheta_\xi)\le C_{\mathfrak d}\log(2N)$.
Given the exponent $\epsilon>0$ in
\eqref{eq3-add-offdiagonal-goal}, use $\epsilon'=\epsilon/2$ in
both counting estimates and
$\log(2N)\le C_\epsilon N^{\epsilon/2}$.
\eqref{eq3-add-resonant-total} and
\eqref{eq3-add-nonresonant-before} then bound $$\sum_{k\in\mathcal K(\xi)}\mathcal S_k(\xi,h)$$
by
\[
 C_{\mathfrak d,\epsilon}N^\epsilon
       \left(Mh+\frac1{Nh}+\sqrt{M/N}\right).
\]
Adding both remainders in \eqref{eq3-add-remainder-sum} proves
\eqref{eq3-add-offdiagonal-goal}. 
%In particular, the remainder
%bounds have not been included implicitly in an estimate for a leading sum.

\medskip
\noindent\textbf{Step 6. The range $0<\tau=Nh<\tau_*$.}
For each $k\in\mathcal K(\xi)$,
\[
 \Theta_k'(j)=-\tau\psi_h(u_k(j))
       \in[-b_\psi\tau,-a_\psi\tau]
       \subset[-1+a_\psi\tau,-a_\psi\tau].
\]
For the last inclusion, use
$(a_\psi+b_\psi)\tau\le2b_\psi\tau_*\le1/2$.
%The distance parameter $a_\psi\tau$ lies in $(0,1/2]$.
%\Cref{lem3-add-derivatives}(i), with integer translation $-1$,
%therefore gives an upper bound $C_{\mathfrak d}\tau^{-1}$ for the
%weighted sum over $j$. Its trivial bound is $C_{\mathfrak d}M$.
%Consequently
%\[
% \left|\sum_{j\in\mathcal J_k}b_k(j)e(\Theta_k(j))\right|
%       \le C_{\mathfrak d}\min\{M,\tau^{-1}\}.
%\]
Applying
\Cref{lem3-add-derivatives}(i),
\eqref{eq3-add-amplitude-variation}, and Abel summation gives
$$ \left|\sum_{j\in\mathcal J_k}b_k(j)e(\Theta_k(j))\right|
\le C_{\mathfrak d}\tau^{-1}.$$
On the other hand, \eqref{eq3-add-amplitude-variation} gives
$$ \left|\sum_{j\in\mathcal J_k}b_k(j)e(\Theta_k(j))\right|
\le |\mathcal J_k|\|b_k\|_{L^\infty([0,2M])}\le C_{\mathfrak d} M$$

Furthermore, \eqref{eq3-add-stationary-count} implies
\[
 \sum_{k\in\mathcal K(\xi)}(N|k|)^{-1/2}
 \le C_{\mathfrak d}\frac{Q_\xi}{\sqrt{NQ_\xi}}
 \le C_{\mathfrak d}\sqrt{M/N}.
\]
Multiplying these bounds and adding
\eqref{eq3-add-remainder-sum} proves
\eqref{eq3-add-tiny-offdiagonal}. 
%All threshold choices preceded the
%choice of $N,M,h,\xi$ and are independent of these varying parameters.
This completes the proof of \Cref{prop3:offdiagonal}.
\end{proof}

\subsection{The diagonal contribution}\label{sec3:diagonal}

The proof of \Cref{prop3:offdiagonal} used $k\ne0$ to
retain a second-derivative lower bound independent of $h$.
For $k=0$, we instead sum over the integer index first and then
use the inverse of $\psi_h$. 
%The constant convention
%$\mathfrak d=(\rho,\theta,\eta,\omega)$ remains in force.
Let $\Omega=|\omega|^2$, extended by zero to $\R$. By
\eqref{eq3-add-mode-expansion},
\begin{equation}
 \mathcal D_h(\xi)=\sum_{m\in\Z}\Omega(m/M)
 \int_{I_0}w_h(u)e\bigl(\xi u-\tau m\psi_h(u)\bigr)\,du,
 \qquad \tau=Nh>0.
 \label{eq3-add-diagonal-original}
\end{equation}
Here and throughout this subsection, $\tau=Nh$. And, we shall use some notations from the previous subsection.
%; it is not an
%additional independent argument of $\mathcal D_h$.

First suppose $0<\tau<\tau_*$. For each fixed $u\in I_0$, the
phase in the integer variable $m$ has slope $-\tau\psi_h(u)$,
which lies in $(-1/2,0)$ at distance at least $a_\psi\tau$ from
$\Z$. The finite geometric progression sum over consecutive relevant integers is
bounded both by $3M$ and by
$2/|1-e(-\tau\psi_h(u))|\le C_{\mathfrak d}/\tau$.
The smooth weight satisfies
\[
 \mathcal V_\R\bigl(\Omega(\cdot/M)\bigr)
 \le\|\Omega\|_\infty+\int_\R|\Omega'(v)|\,dv,
\]
independently of $M$. Abel summation, followed by
$\|w_h\|_{L^1(\R)}\le\|\eta\|_{L^2(\R)}^2$, therefore gives
\begin{equation}
 \sup_{\xi\in\R}|\mathcal D_h(\xi)|
       \lesssim_{\mathfrak d}\min\{M,\tau^{-1}\}.
 \label{eq3-add-tiny-diagonal}
\end{equation}
The integral of $\min\{M,\tau^{-1}\}$ on $(0,\tau_*)$ is
$M\tau_*\le1$ if $M\tau_*\le1$. Otherwise it equals
\[
 \int_0^{1/M}M\,d\tau+\int_{1/M}^{\tau_*}\frac{d\tau}{\tau}
 =1+\log(M\tau_*)\le1+\log M.
\]
Together with \eqref{eq3-add-tiny-offdiagonal} and
$\sqrt{M/N}\le1$, this proves
\begin{equation}
 \frac1L\int_0^{\tau_*}\sup_{\xi\in\R}
          |\mathcal A_{\tau/N}(\xi)|\,d\tau
 \lesssim_{\mathfrak d}\frac{\log(2M)}L+MN^{-3/2}.
 \label{eq3-add-tiny-integral}
\end{equation}
The use of \eqref{eq3-add-tiny-integral} requires $N\ge N_{\rm off}$.
%and $L\ge\tau_*$. 
%The point $h=0$ has Lebesgue measure zero and
%is not assigned a bound involving $h^{-1}$.

\begin{prop}\label[prop]{prop3:diagonal}
There is $h_0=h_0(\mathfrak d)\in(0,1]$ such that, for every integer
$N\ge2$, real $M\in[1,N^\theta]$, and
$0<h\le\min\{H,h_0\}$,
\begin{equation}
 \sup_{\xi\in\R}|\mathcal D_h(\xi)|
       \lesssim_{\mathfrak d}1+\tau^{-1},\qquad \tau=Nh.
 \label{eq3-add-diagonal-absolute}
\end{equation}
For each fixed integer $d\ge3$, there is
$h_d=h_d(\mathfrak d,d)\in(0,h_0]$ such that, if also
$h\le h_d$ and $\tau\ge1$, then
\begin{equation}
 \sup_{\xi\in\R}|\mathcal D_h(\xi)|
 \lesssim_{\mathfrak d,d}\tau^{-1}
       +\left(\frac M{\tau^{d-1}}\right)^{1/(2^d-2)}
       +(M\tau)^{-1}.
 \label{eq3-add-diagonal-high-derivative}
\end{equation}
%In particular, neither the constants nor $h_0,h_d$ depend on
%$N,M,h,\xi$ within these ranges.
\end{prop}
\begin{proof}
The proof is divided into three steps.

\medskip
\noindent\textbf{Step 1. An inverse on a common open interval.}
Retain the intervals $I_0,I_1,I_2,J_\psi$ fixed in the proof of
\Cref{prop3:offdiagonal}, and put
\[
 J=J_\psi,\qquad J^+=[r_-/2,2r_+],\qquad
 \mathcal O_{\rm inv}=(r_-/4,4r_+).
\]
The last interval is open and contains the closed interval $J^+$.
For $u\in I_1$ and $0\le h\le1$,
\[
 \psi_h'(u)=\int_0^1\psi''(u+sh)\,ds<0,\qquad
 |\psi_h'(u)|\ge\min_{I_2}|\psi''|=:m_2>0.
\]
For each integer $s\ge0,0\le h\le 1$, differentiation under the integral in
\eqref{eq3-add-divided-difference} gives
\begin{equation}\label{eq3-3-0}
	\sup_{u\in I_1}|\psi_h^{(s)}(u)-\psi^{(s+1)}(u)|
	\le\frac h2\max_{I_2}|\psi^{(s+2)}|
	\le C_{\mathfrak d,s}h.
\end{equation}
%\[
% \sup_{u\in I_1}|\psi_h^{(s)}(u)-\psi^{(s+1)}(u)|
% \le\frac h2\max_{I_2}|\psi^{(s+2)}|
% \le C_{\mathfrak d,s}h.
%\]
%At $h=0$, the values at $a_1,b_1$ surround the closure of
%$\mathcal O_{\rm inv}$ with strictly positive margins.
Note that $\psi'(I_0)\subset\operatorname{int}J$.
Choose a closed interval $J_0\subset\operatorname{int}J$ containing
$\psi'(I_0)$ in its interior. The uniform estimate \eqref{eq3-3-0}, with $s=0$,
permits a single $h_0\in(0,1]$ such that
\[
 \mathcal O_{\rm inv}\subset\psi_h((a_1,b_1)),\qquad
 \psi_h(I_0)\subset\operatorname{int}J_0
                                \quad(0\le h\le h_0).
\]
Both $J_0$ and $h_0$ are chosen from $\rho,I_0,I_1,I_2$ and hence
from $\mathfrak d$. For $0\le h\le h_0$, the strict monotonicity of $\psi_h$ and the inverse function
theorem define a smooth real function
\[
 \Psi_h:\mathcal O_{\rm inv}\to(a_1,b_1),\qquad
 \psi_h(\Psi_h(v))=v.
\]
%This definition includes $h=0$.

We verify uniformity of all inverse derivatives. For $v\in J^+$,
write $x_h(v)=\Psi_h(v)$ and $x_0(v)=\Psi_0(v)$. The mean value
theorem applied to $\psi_0=\psi'$ gives
\[
 m_0|x_h(v)-x_0(v)|
 \le|\psi_0(x_h(v))-\psi_0(x_0(v))|
 =|\psi_0(x_h(v))-\psi_h(x_h(v))|
 \le\|\psi_0-\psi_h\|_{C^0(I_1)}.
\]
Thus $\Psi_h\to\Psi_0$ uniformly on $J^+$ by \eqref{eq3-3-0}.
On the open interval $(0,\infty)$, for $0\le h\le h_0$, define
\[
 R_{1,h}=1/\psi_h',\qquad R_{s+1,h}=R_{s,h}'/\psi_h'\quad(s\ge1).
\]
Differentiating the inverse identity proves by induction that
$\Psi_h^{(s)}=R_{s,h}\circ\Psi_h$ on $\mathcal O_{\rm inv}$.
%For example,
%\[
% R_{2,h}=-\frac{\psi_h''}{(\psi_h')^3},\qquad
% R_{3,h}=\frac{3(\psi_h'')^2}{(\psi_h')^5}
%                         -\frac{\psi_h^{(3)}}{(\psi_h')^4}.
%\]

%At every fixed order the recurrence uses finitely many derivatives of
%$\psi_h$ and reciprocal powers of $\psi_h'$. Their uniform bounds,
%the lower bound $m_2$, and $C^s$ convergence of $\psi_h$ prove
%uniform bounds and convergence of each $R_{s,h}$ on $I_1$.
%In the composition, use
%\[
% \begin{split}
% |R_{s,h}(x_h(v))-R_{s,0}(x_0(v))|
% \le{}&\|R_{s,h}-R_{s,0}\|_{C^0(I_1)}\\
% &+|R_{s,0}(x_h(v))-R_{s,0}(x_0(v))|.
% \end{split}
%\]
%Uniform continuity of the fixed function $R_{s,0}$ then proves
%\[
% \sup_{0\le h\le h_0}\|\Psi_h\|_{C^s(J^+)}\le C_{\mathfrak d,s},
% \qquad \|\Psi_h-\Psi_0\|_{C^s(J^+)}\longrightarrow0
%                                  \quad(h\downarrow0)
%\]
%for every integer $s\ge0$. Endpoint derivatives are restrictions
%from the open inverse domain.

For each fixed integer $s\ge1$, the recurrence defining $R_{s,h}$
involves only $\psi_h',\ldots,\psi_h^{(s)}$ and divisions by
powers of $\psi_h'$. Since $|\psi_h'|\ge m_2>0$ on $I_1$,
the uniform derivative bounds and the convergence
$\psi_h\to\psi_0$ in $C^s(I_1)$ imply
\[
\sup_{0\le h\le h_0}\|R_{s,h}\|_{C^0(I_1)}
\le C_{\mathfrak d,s},
\qquad
\|R_{s,h}-R_{s,0}\|_{C^0(I_1)}\longrightarrow0
\quad(h\downarrow0).
\]
Recall that
$\Psi_h^{(s)}=R_{s,h}\circ\Psi_h$ and
$\Psi_h(J^+)\subset I_1$. Hence
\[
\|\Psi_h^{(s)}\|_{C^0(J^+)}
\le \|R_{s,h}\|_{C^0(I_1)}
\le C_{\mathfrak d,s}.
\]
To prove convergence, use the mean value theorem for the
smooth function $R_{s,0}$ on $I_1$:
\[
\begin{aligned}
\|\Psi_h^{(s)}-\Psi_0^{(s)}\|_{C^0(J^+)}
\le \|R_{s,h}-R_{s,0}\|_{C^0(I_1)}
+
\|R_{s,0}'\|_{C^0(I_1)}
\|\Psi_h-\Psi_0\|_{C^0(J^+)}.
\end{aligned}
\]
Both terms tend to zero as $h\downarrow0$. 
%The first does so by the preceding
%estimate, and the second by the uniform convergence of the
%inverses already proved.
Applying these bounds at orders $1,\ldots,s$, together with the
order-zero bounds, gives
\[
\sup_{0\le h\le h_0}\|\Psi_h\|_{C^s(J^+)}
\le C_{\mathfrak d,s},
\qquad
\|\Psi_h-\Psi_0\|_{C^s(J^+)}\longrightarrow0
\quad(h\downarrow0)
\]
for every integer $s\ge0$. 
%Here derivatives are taken on
%$\mathcal O_{\rm inv}$ and restricted to the closed interval $J^+$.

Define
\begin{equation}
 a_h(v)=
 \begin{cases}
 w_h(\Psi_h(v))/|\psi_h'(\Psi_h(v))|,&v\in\mathcal O_{\rm inv},\\
 0,&v\notin\mathcal O_{\rm inv}.
 \end{cases}
 \label{eq3-add-inverse-amplitude}
\end{equation}
A nonzero numerator requires $\Psi_h(v)\in I_0$ and hence
$v\in\psi_h(I_0)\subset\operatorname{int}J_0$.
Thus $a_h$ extends smoothly by zero and
$$\sup_{0\le h\le h_0}\|a_h\|_{C^s(\R)}\le C_{\mathfrak d,s}$$ for every fixed $s$.

%Choose, once from $J,J^+$, a real
%$\chi_{\rm inv}\in C_c^\infty(\operatorname{int}J^+)$ equal to one
%on an open neighborhood of $J$. Multiplying $\Psi_h$ by
%$\chi_{\rm inv}$ on $\mathcal O_{\rm inv}$ and extending by zero
%defines $\Psi_h^{\rm ext}\in C_c^\infty(\R;\R)$.
%Its fixed-order norms are bounded by $C_{\mathfrak d,s}$; only its
%restriction near $J$ is asserted to be an inverse.

Choose a real-valued function $\chi_{\rm inv}\in C_c^\infty(\R)$,
depending only on $J$ and $J^+$, such that
\[
\operatorname{supp}\chi_{\rm inv}\subset\operatorname{int}J^+,
\qquad
\chi_{\rm inv}=1 \quad\text{on an open neighborhood of }J.
\]
For $0\le h\le h_0$, define
\[
\Psi_h^{\rm ext}(v)=
\begin{cases}
\chi_{\rm inv}(v)\Psi_h(v),&v\in\mathcal O_{\rm inv},\\
0,&v\notin\mathcal O_{\rm inv}.
\end{cases}
\]
Since $\operatorname{supp}\chi_{\rm inv}$ is compactly contained
in $\mathcal O_{\rm inv}$, this defines a smooth function on
$\R$ with compact support. The product rule and the uniform
derivative bounds for $\Psi_h$ on $J^+$ give
\[
\sup_{0\le h\le h_0}
\|\Psi_h^{\rm ext}\|_{C^s(\R)}
\le C_{\mathfrak d,s}
\qquad(s\ge 0).
\]
%On the neighborhood of $J$ where $\chi_{\rm inv}=1$, we have
%$\Psi_h^{\rm ext}=\Psi_h$, and hence
%$\psi_h(\Psi_h^{\rm ext}(v))=v$.
%No inverse identity is asserted outside that neighborhood.

In \eqref{eq3-add-diagonal-original}, the decreasing substitution
$v=\psi_h(u)$ yields
\begin{equation}
 \mathcal D_h(\xi)=\sum_{m\in\Z}\Omega(m/M)
 \int_\R a_h(v)e\bigl(\xi\Psi_h^{\rm ext}(v)-\tau mv\bigr)\,dv.
 \label{eq3-add-diagonal-inverse}
\end{equation}
%Only $v\in J_0$ contributes, so the extensions do not alter the
%integral.

\medskip
\noindent\textbf{Step 2. Poisson summation and the relevant frequencies.}
Apply \Cref{lem3-add-dual-formula} with
$a=a_h$ and $\Psi=\Psi_h^{\rm ext}$.
Their supports and derivative upper bounds are supplied by the estimates
above. Also on $J$, 
\[
 |(\Psi_h^{\rm ext})'(v)|=|\Psi_h'(v)|
 =\frac1{|\psi_h'(\Psi_h(v))|}
 \ge\left(\max_{I_2}|\psi''|\right)^{-1}>0.
\]
Thus the fixed data in \Cref{lem3-add-dual-formula} depend only on $\mathfrak d$, and
\eqref{eqB:absolute} proves
\eqref{eq3-add-diagonal-absolute}.
For $\tau\ge1$ and $|\xi|\le C_\infty M\tau$, where
$C_\infty=C_\infty(\mathfrak d)$ is chosen as in
\Cref{lem3-add-dual-formula}, \eqref{eqB:poisson-formula}
gives
\begin{equation}
 \begin{split}
 \mathcal D_h(\xi)=\frac1\tau
 \sum_{\ell\in\tau J\cap\Z}a_h(\ell/\tau)
 \Omega\left(\frac{\xi\Psi_h'(\ell/\tau)}{M\tau}\right)
 e\bigl(\xi\Psi_h(\ell/\tau)\bigr)
 +O_{\mathfrak d}((M\tau)^{-1}).
 \end{split}
 \label{eq3-add-dual-formula}
\end{equation}
%Here $\ell$ is the integer dual summation index; it is unrelated to
%$r_k=\xi/(Nk)$ in \eqref{eq3-add-normalized-phase}.
%The support of $a_h$ is contained in $J_0\subset J$.
There are fixed $0<a_{\rm inv}\le b_{\rm inv}$, depending only on
$\mathfrak d$, such that
$a_{\rm inv}\le|\Psi_h'|\le b_{\rm inv}$ on $J$.
Since $\Psi_h'<0$ and $\operatorname{supp}\Omega\subset(1/2,2)$,
a nonzero leading sum in \eqref{eq3-add-dual-formula} requires
\[
 \xi<0,\qquad
 \frac1{2b_{\rm inv}}M\tau<|\xi|<\frac2{a_{\rm inv}}M\tau.
\]
For $\tau\ge1$ and $|\xi|>C_\infty M\tau$, \eqref{eqB:large-frequency}
gives $|\mathcal D_h(\xi)|\le C_{\mathfrak d}(M\tau)^{-1}$.
%For frequencies inside this range but with zero leading sum, the
%remainder in \eqref{eq3-add-dual-formula} gives the same bound.

\medskip
\noindent\textbf{Step 3. Derivative estimates for the dual sum.}
On $\mathcal O_{\rm inv}$, the limiting inverse is
\begin{equation}
 \Psi_0(v)=(v/\rho)^{1/(\rho-1)}.
 \label{eq3-add-limiting-inverse}
\end{equation}
Write $a_\rho=1/(\rho-1)<0$. For every integer $s\ge1$,
\[
 \Psi_0^{(s)}(v)
 =\rho^{-a_\rho}a_\rho(a_\rho-1)\cdots(a_\rho-s+1)v^{a_\rho-s}.
\]
Its sign is $(-1)^s$ and
$b_s:=\min_J|\Psi_0^{(s)}|>0$. Also set
$B_s:=\max_J|\Psi_0^{(s)}|$.
Fix $d\ge3$. Convergence in $C^d(J^+)$ permits one
$h_d=h_d(\mathfrak d,d)\in(0,h_0]$ such that, for $0\le h\le h_d$,
$\sup_J|\Psi_h^{(s)}-\Psi_0^{(s)}|\le b_s/2$ simultaneously for
$1\le s\le d$. Hence,
\begin{equation}
 c_s\le|\Psi_h^{(s)}(v)|\le C_s
       \quad(v\in J,\ 0\le h\le h_d,\ 1\le s\le d),
 \quad c_s=b_s/2,\quad C_s=B_s+b_s/2.
 \label{eq3-add-inverse-derivatives}
\end{equation}
These constants depend only on $\rho,J,s$. Each derivative retains sign $(-1)^s$.

Suppose $\tau\ge2$ and the leading sum in
\eqref{eq3-add-dual-formula} is not identically zero.
Set $\zeta_{M,\tau,\xi}=\xi/(M\tau)$. The frequency restrictions
above place this real parameter in a fixed compact subset of
$(-\infty,0)$, separated from zero.
For such a parameter $\zeta$, define the smooth function
\[
 b_{h,\zeta}(v)=a_h(v)
       \Omega\bigl(\zeta(\Psi_h^{\rm ext})'(v)\bigr),\qquad v\in\R,0\le h\le h_d.
\]
Its support is contained in $J_0$, and the product rule gives
\[
 b_{h,\zeta}'=a_h'\Omega\bigl(\zeta(\Psi_h^{\rm ext})'\bigr)
 +a_h\Omega'\bigl(\zeta(\Psi_h^{\rm ext})'\bigr)
                                  \zeta(\Psi_h^{\rm ext})''.
\]
Every factor is uniformly bounded using $\mathfrak d$. Therefore
\[
 \mathcal V_\R(b_{h,\zeta})
 \le\|b_{h,\zeta}\|_{L^{\infty}(\R)}+|J_0|\sup_\R|b_{h,\zeta}'|
 \le C_{\mathfrak d}.
\]
%The weight $x\mapsto b_{h,\zeta}(x/\tau)$ on $\tau J$ has the same
%bound, since
%\[
% \int_{\tau J}\tau^{-1}|b_{h,\zeta}'(x/\tau)|\,dx
% =\int_J|b_{h,\zeta}'(v)|\,dv.
%\]
For $\tau>0$, the substitution $x=\tau v$ gives that 
\[
\begin{aligned}
\|b_{h,\zeta}(\cdot/\tau)\|_{L^\infty(\tau J)}
&=\|b_{h,\zeta}\|_{L^\infty(J)},\\
\int_{\tau J}
\left|\frac{d}{dx}b_{h,\zeta}(x/\tau)\right|\,dx
&=\frac1\tau\int_{\tau J}|b_{h,\zeta}'(x/\tau)|\,dx
=\int_J|b_{h,\zeta}'(v)|\,dv.
\end{aligned}
\]
Consequently,
\begin{equation}\label{eq3-3-1}
	\mathcal V_{\tau J}\bigl(b_{h,\zeta}(\cdot/\tau)\bigr)
	=\mathcal V_J(b_{h,\zeta})
	\le C_{\mathfrak d}.
\end{equation}
%with a constant independent of $\tau$ and of the admissible
%parameters $h,\zeta$.

Based on the form of the leading sum in \eqref{eq3-add-dual-formula}, we define the real phase by
\[
 q_{h,\tau,\xi}:\tau\mathcal O_{\rm inv}\to\R,
 \qquad q_{h,\tau,\xi}(x)=\xi\Psi_h(x/\tau),\qquad 0\le h\le h_d.
\]
%Its open domain contains the closed summation interval $\tau J$.
\eqref{eq3-add-inverse-derivatives} implies that 
\[
 q_{h,\tau,\xi}^{(s)}(x)=\xi\tau^{-s}\Psi_h^{(s)}(x/\tau),\qquad
 |q_{h,\tau,\xi}^{(s)}(x)|\asymp_{\mathfrak d,s}|\xi|\tau^{-s}
                    \quad(1\le s\le d,\ x\in\tau J).
\]

We apply \Cref{lem3-add-derivatives}(iii) with
$L=\tau$, $\mathcal F=|\xi|$,
$a_*=r_-$, and $b_*=r_+$. 
%All required derivative comparisons, not just
%the comparison at order $d$, have been verified.
Applying
\Cref{lem3-add-derivatives}(iii),
\eqref{eq3-3-1}, and Abel summation gives
\[
 \begin{split}
 &\frac1\tau\left|\sum_{\ell\in\tau J\cap\Z}
 b_{h,\zeta_{M,\tau,\xi}}(\ell/\tau)
                         e(q_{h,\tau,\xi}(\ell))\right|
 \lesssim_{\mathfrak d,d}
   \tau^{-1}+(|\xi|\tau^{-d})^{1/(2^d-2)}+|\xi|^{-1}\\
 &\hspace{8cm}\lesssim_{\mathfrak d,d}
   \tau^{-1}+(M/\tau^{d-1})^{1/(2^d-2)}+(M\tau)^{-1}.
 \end{split}
\]

The second inequality uses $|\xi|\asymp_{\mathfrak d}M\tau$.
The amplitude is zero outside $\tau J_0$, so it is the leading
sum in \eqref{eq3-add-dual-formula}. Adding its remainder proves
\eqref{eq3-add-diagonal-high-derivative} for the present frequencies.
The complementary frequencies were already bounded by
$C_{\mathfrak d}(M\tau)^{-1}$.
For $1\le\tau<2$, \eqref{eq3-add-diagonal-absolute} gives
at most $2C_{\mathfrak d}$.
Then by $\tau^{-1}>1/2$, increasing $C_{\mathfrak d,d}$ covers this interval and proves
\Cref{prop3:diagonal}.
\end{proof}

\subsection{Completion of the proof of \texorpdfstring{\Cref{thm3-1}}{Theorem~\getrefnumber{thm3-1}}}\label{sec3-13}

Take $d=d_\theta$ and $P=2^d\ge 16$, as in
\eqref{eq3-add-parameters}. These integers are fixed by $\theta$.
Choose an integer $N_0$ at least
\[
 \max\left\{N_{\rm off},\,
 \left\lceil h_d^{-1/(1-\sigma)}\right\rceil,\,
 \left\lceil 2^{2/\sigma}\right\rceil\right\}.
\]
Then, for every $N\ge N_0>1$, $H\le h_d$ and
$$L>N^{\sigma/2}\ge 2\ge 1\ge \tau_{*}>0.$$
%The number $N_0$ depends only on $\mathfrak d$; it is uniform for $1\le M\le N^\theta$.

We estimate the normalized scalar integral in
\eqref{eq3-add-dual-time-reduction}. First,
\eqref{eq3-add-tiny-integral} gives
\begin{equation}
 \frac1L\int_0^{\tau_*}\sup_{\xi\in\R}
       |\mathcal A_{\tau/N}(\xi)|\,d\tau
 \lesssim_{\mathfrak d}N^{-\sigma}\log(2N)+MN^{-3/2}
 \lesssim_{\mathfrak d}N^{-\sigma/2}.
 \label{eq3-add-final-tiny}
\end{equation}
%For the second inequality, use
%$\log(2N)\le C_\sigma N^{\sigma/2}$ and
%\[
% MN^{-3/2}\le N^{\theta-3/2},\qquad
% \theta-3/2+\sigma/2=(3\theta-5)/4<0.
%\]
%The constant $C_\sigma$ is included in $C_{\mathfrak d}$ because
%$\sigma=(1-\theta)/2$.

For the off-diagonal contribution on $[\tau_*,L]$, substitute
$h=\tau/N$ into \eqref{eq3-add-offdiagonal-goal}. The three terms
have the respective normalized integrals:
\[
 \begin{split}
 \frac1L\int_{\tau_*}^L Mh\,d\tau
   \le\frac{ML}{2N},\qquad
 \frac1L\int_{\tau_*}^L\frac{d\tau}{Nh}
   =\frac{\log(L/\tau_*)}{L},\qquad
 \frac1L\int_{\tau_*}^L\sqrt{M/N}\,d\tau\le\sqrt{M/N}.
 \end{split}
\]
Consequently, for every fixed $\epsilon>0$,
\begin{equation}
 \begin{split}
 \frac1L\int_{\tau_*}^L\sup_{\xi\in\R}
                  |\mathcal O_{\tau/N}(\xi)|\,d\tau
 \lesssim_{\mathfrak d,\epsilon}N^\epsilon
       \left(\frac{ML}{N}+\frac{\log(2L)}L+\sqrt{\frac MN}\right)
 \lesssim_{\mathfrak d,\epsilon}
                          N^{-\sigma+\epsilon}\log(2N).
 \end{split}
 \label{eq3-add-final-offdiagonal}
\end{equation}
%The three scale relations used here are
%\[
% ML/N\le N^{\theta+\sigma-1}=N^{-\sigma},\qquad
% L^{-1}=N^{-\sigma},\qquad
% \sqrt{M/N}\le N^{(\theta-1)/2}=N^{-\sigma}.
%\]
Choose $\epsilon=\sigma/4$.
Then the quantity in \eqref{eq3-add-final-offdiagonal} is at most
$C_{\mathfrak d}N^{-\sigma/2}$.

On $[\tau_*,N^{\sigma/2}]$,
\eqref{eq3-add-diagonal-absolute} is at most
$C_{\mathfrak d}(1+\tau_*^{-1})$. Hence,
\begin{equation}
 \frac1L\int_{\tau_*}^{N^{\sigma/2}}\sup_{\xi\in\R}
               |\mathcal D_{\tau/N}(\xi)|\,d\tau
 \lesssim_{\mathfrak d}N^{-\sigma/2}.
 \label{eq3-add-final-moderate}
\end{equation}
For $1<N^{\sigma/2}\le\tau\le L$, the definition of $d$ gives
$(d-1)\sigma/2\ge1$. Thus
\[
 \frac M{\tau^{d-1}}
 \le N^{\theta-(d-1)\sigma/2}\le N^{\theta-1}=N^{-2\sigma}.
\]
At present, all assumptions of \eqref{eq3-add-diagonal-high-derivative} hold,
including $h=\tau/N\le H\le h_d$ and $\tau\ge1$.
Its three terms satisfy
\begin{align*}
 \tau^{-1}&\le N^{-\sigma/2}\le N^{-\sigma/P},\\
 (M/\tau^{d-1})^{1/(P-2)}
   &\le N^{-2\sigma/(P-2)}\le N^{-\sigma/P},\\
 (M\tau)^{-1}&\le\tau^{-1}\le N^{-\sigma/P}.
\end{align*}
%Here $M\ge1$, $P>2$, and $2/(P-2)\ge1/P$.
%Dividing by $L$ and using an integration interval of length at most
%$L$ proves
By the above, we have 
\begin{equation}
 \frac1L\int_{N^{\sigma/2}}^L\sup_{\xi\in\R}
                 |\mathcal D_{\tau/N}(\xi)|\,d\tau
 \lesssim_{\mathfrak d}N^{-\sigma/P}.
 \label{eq3-add-final-large}
\end{equation}
%The constant $C_{\mathfrak d,d}$ in
%\eqref{eq3-add-diagonal-high-derivative} is now a function of
%$\mathfrak d$ alone because $d=d_\theta$.

Use $$\sup_\xi|\mathcal A_h(\xi)|
\le\sup_\xi|\mathcal O_h(\xi)|+\sup_\xi|\mathcal D_h(\xi)|$$
only on $[\tau_*,L]$, and use
\eqref{eq3-add-final-tiny} on $(0,\tau_*)$.
\eqref{eq3-add-final-offdiagonal},
\eqref{eq3-add-final-moderate}, and \eqref{eq3-add-final-large}
then give \eqref{eq3-add-main-bound} for $N\ge N_0$ by
\eqref{eq3-add-dual-time-reduction}.
%The length $|I_0|$ in that reduction is fixed by $\eta$ and is
%included in $C_{\mathfrak d}$.

For $2\le N<N_0$, there are at most $2N_0^\theta+1$ integers
with $\omega(m/M)\ne0$. The triangle inequality gives
\[
 \left\|\int_\R V(u)\,du\right\|_2
 \lesssim \|\eta\|_{L^1(\R)}\|\omega\|_{L^\infty(\R)}N_0^\theta\lesssim_{\mathfrak d} N_{0}^{\theta+\sigma/P}N_{0}^{-\sigma/P}\lesssim_{\mathfrak d} N^{-\sigma/P}.
\]
%Since $B_{\mathfrak d}^2
%\le B_{\mathfrak d}^2N_0^{\sigma/P}N^{-\sigma/P}$ on this
%remaining range, increasing $C_{\mathfrak d}$ by this fixed amount
%proves \eqref{eq3-add-main-bound} for every $N\ge2$.
This completes the proof of \Cref{thm3-1}.\hfill$\square$

%\clearpage
\appendix
\section{Analytic estimates}\label[appendix]{app:derivatives}

We state some estimates for exponential sums used in \Cref{sec3:main}. By Abel summation, these estimates also work for weights with bounded variation norms.
%Dependencies in this appendix refer only to the fixed data of the
%individual assertion. They do not include the data $\mathfrak d$
%from \Cref{sec3:offdiagonal}. 
%Numerical constants denoted
%$C_{\rm abs}$ are absolute.

\begin{lemma}\label[lemma]{lem3-add-derivatives}
Let $I=[a_-,a_+]$, where $a_-<a_+$ are real, and let
$\mathcal O\subset\R$ be an open interval containing $I$.
The real phase $q:\mathcal O\to\R$ has the regularity specified
in the following assertions. 
%All sums are over integer indices.
\begin{itemize}
\item[(i)] Suppose $q\in C^1(\mathcal O;\R)$ and $q'$ is
monotone on $I$. If $q'(I)\subset[r+\delta,r+1-\delta]$ for
some $r\in\Z$ and $0<\delta\le1/2$, then
\[
 \left|\sum_{n\in I\cap\Z}e(q(n))\right|
                         \lesssim \delta^{-1}.
\]
\item[(ii)] Suppose $q\in C^2(\mathcal O;\R)$,
$\lambda\le|q''(x)|\le C_0\lambda$ for every $x\in I$, where
$\lambda>0$, $C_0\ge1$, and $|I|\le L$ with $L\ge1$. Then
\[
 \left|\sum_{n\in I\cap\Z}e(q(n))\right|
          \lesssim_{C_0}\bigl(L\lambda^{1/2}+\lambda^{-1/2}\bigr).
\]
\item[(iii)] Fix $0<a_*<b_*$, an integer $d\ge3$, and constants
$0<c_s\le C_s$ for $1\le s\le d$; denote these data by
$\mathfrak a_d=(a_*,b_*,d,(c_s)_{s=1}^d,(C_s)_{s=1}^d)$.
Suppose $L\ge2$, $\mathcal F>0$, $I\subset[a_*L,b_*L]$,
and $q\in C^d(\mathcal O;\R)$ satisfies
\[
 c_s\mathcal F L^{-s}\le|q^{(s)}(x)|\le C_s\mathcal F L^{-s}
                         \quad(x\in I,\ 1\le s\le d).
\]
Then, with $P_d=2^d$,
\begin{equation}
 \frac1L\left|\sum_{n\in I\cap\Z}e(q(n))\right|
 \lesssim_{\mathfrak a_d}
 L^{-1}+(\mathcal F L^{-d})^{1/(P_d-2)}+\mathcal F^{-1}.
 \label{eq3-add-derivative-test}
\end{equation}
\end{itemize}
The implicit constants in (i), (ii), and (iii) are independent of
$\mathcal O,I,q,r,\delta,\lambda,L,\mathcal F$.
\footnote
{Every estimate holds with the original scale when the sum is
restricted to a closed subinterval of $I$, including a singleton.}
%An empty sum is zero. For a pointwise-defined function $a:I\to\C$
%of bounded variation, the corresponding weighted sum satisfies the
%same bound multiplied by $\mathcal V_I(a)$, with no additional
%constant depending on $a$.
\end{lemma}
\begin{proof}
	
\medskip	
\noindent\textbf{Part I. First derivative. }
Empty and one-term sums satisfy the claimed bound. Otherwise, let
$n_-=\min(I\cap\Z)<n_+=\max(I\cap \Z)$. Replacing $q(t)$ by $q(t)-rt$ does not
change $e(q(n))$ at integer $n$. Thus we may assume
$\delta\le q'\le1-\delta$ on $I$.
For $n_-\le n<n_+$, define
\[
 t_n=q(n+1)-q(n)=\int_0^1q'(n+s)\,ds,\qquad z_n=e(q(n)).
\]
The sequence $t_n$ is monotone and lies in $[\delta,1-\delta]$.
For $0<t<1$, let
\[
 B(t)=(1-e(t))^{-1}=\frac12+\frac i2\cot(\pi t).
\]
The identity $z_n=B(t_n)(z_n-z_{n+1})$, summed over
$n_-\le n<n_+$, and the remaining term $z_{n_+}$ give
\[
 \begin{split}
 \sum_{n=n_-}^{n_+}z_n={}&B(t_{n_-})z_{n_-}
       +(1-B(t_{n_+-1}))z_{n_+}
 +\sum_{n=n_-+1}^{n_+-1}
       (B(t_n)-B(t_{n-1}))z_n.
 \end{split}
\]
Since $|e(q(n))|=1$, it is enough to bound the coefficients here.
On $[\delta,1-\delta]$,
\[
 \sup|B|\le\frac1{2\sin(\pi\delta)}\lesssim \delta^{-1},
 \qquad
 \int_\delta^{1-\delta}|B'(t)|\,dt
       =\cot(\pi\delta)\lesssim \delta^{-1}.
\]
%At $\delta=1/2$ the integral equals zero. 
Monotonicity of $t_n$ gives
$$\left| \sum_{n=n_-}^{n_+}z_n\right|\lesssim \delta^{-1}+\int_\delta^{1-\delta}|B'(t)|\,dt\lesssim \delta^{-1}.$$
%bounds the sum of absolute increments of $B(t_n)$ by this total
%variation. The endpoint coefficient $1-B(t_{n_+-1})$ is bounded by
%$1+\sup|B|$. 
This proves \Cref{lem3-add-derivatives}(i).
%with an absolute constant independent of the integer $r$.

\medskip
\noindent\textbf{Part II. Second derivative.}
%We use \cite[equation~(1), $k=2$, author's version]{HeathBrown2017}.
Empty and one-term sums satisfy the claimed bound. Otherwise, let
$n_-:=\min(I\cap\Z)<n_+:=\max(I\cap \Z)$. Put $1\le D:=n_+-n_-\le L$.
The translated phase $x\mapsto q(n_-+x)$ is defined on a suitable open
neighborhood of $[0,D]$ and has the same derivative bounds on $[0,D]$.
The sign of $q''$ is constant, since it is continuous and nonzero.
By the case $k=2$ of \cite[Equation~(1)]{HeathBrown2017}, we have 
\[
\left|\sum_{n=1}^{D}e(q(n+n_{-}))\right|
\lesssim_{C_0}\bigl(D\sqrt{\lambda}+\lambda^{-1/2}\bigr).
\]
%\[
% C_{C_0}\bigl(D\sqrt\lambda+\lambda^{-1/2}\bigr).
%\]
Adding the index $n=0$ costs at most one. This can be absorbed by
$L\sqrt\lambda+\lambda^{-1/2}$. 
This proves \Cref{lem3-add-derivatives}(ii).
%The same inequality also handles a one-point summation interval.
%Thus the constant depends on $C_0$ alone, not on the translated
%interval or its endpoints.

\medskip
\noindent\textbf{Part III. Higher derivatives.}
%Apply \cite[Theorem~1.2]{LST2025} with integer
%$l=d-2$, $M_0=a_*L$, and a fixed dilation constant
%$C_{\rm dil}>\max\{1,b_*/a_*\}$.
%For a nondegenerate closed subinterval $[b_-,b_+]\subseteq I$,
%its half-open version lies in $[M_0,C_{\rm dil}M_0)$.
%When $M_0>2$, the derivative assumptions give, for $1\le s\le d$,
%\[
% |q^{(s)}(x)|\asymp_{a_*,s,c_s,C_s}
%                         \mathcal F M_0^{-s}.
%\]
%The A-process hypotheses of \cite[Theorem~1.2]{LST2025} are
%therefore satisfied at every required order. With $Q_l=2^l$, its
%bound is
%\[
% C_{\mathfrak a_d}
% \left(\mathcal F^{1/(4Q_l-2)}
%             M_0^{1-d/(4Q_l-2)}+\mathcal F^{-1}M_0\right).
%\]
%The relation $4Q_l-2=2^d-2=P_d-2$ and division by $L$ turn this
%into
%\[
% C_{\mathfrak a_d}
% \left((\mathcal F L^{-d})^{1/(P_d-2)}+\mathcal F^{-1}\right).
%\]
%Restoring a closed right endpoint adds at most $L^{-1}$.
%A singleton is covered by the same endpoint term.
%If $M_0\le2$, then $L\le2/a_*$ and
%\[
% \#(I\cap\Z)\le(b_*-a_*)L+1
%                    \le2(b_*-a_*)/a_*+1.
%\]
%Thus $C_{a_*,b_*}L^{-1}$ bounds the normalized sum in this case.
%This proves \eqref{eq3-add-derivative-test}, also on every closed
%subinterval with the original scale. The argument uses bounds for
%all derivatives of orders $1,\ldots,d$; it is not an application
%under a hypothesis on $q^{(d)}$ alone.
Set
$M_0=a_*L$. If $M_0\le2$, then
\[
\begin{aligned}
\frac1L\left|\sum_{n\in I\cap\Z}e(q(n))\right|
\le \frac{(b_*-a_*)L+1}{L}
\le \left(1+\frac{2(b_*-a_*)}{a_*}\right)L^{-1}.
\end{aligned}
\]
A singleton also contributes at most $L^{-1}$, and an empty sum
is zero. We may therefore assume that $M_0>2$ and $|[a_{-},a_{+}]\cap \Z|>1$.

Choose $C_{\rm dil}=1+b_*/a_*$. Then
$
[a_-,a_+)\subset[M_0,C_{\rm dil}M_0).
$
Since $L^{-s}=a_*^sM_0^{-s}$, the derivative assumptions imply
\[
c_sa_*^s\mathcal F M_0^{-s}
\le |q^{(s)}(x)|
\le C_sa_*^s\mathcal F M_0^{-s}
\qquad(x\in I,\ 1\le s\le d).
\]
Thus \cite[Theorem~1.2]{LST2025} applies to the sum over
$a_-\le n<a_+$ with $l=d-2$. It gives
\[
\left|\sum_{a_-\le n<a_+}e(q(n))\right|
\lesssim_{\mathfrak a_d}
\left(
\mathcal F^{1/(P_d-2)}M_0^{\,1-d/(P_d-2)}
+\mathcal F^{-1}M_0
\right).
\]
Substituting $M_0=a_*L$ and dividing by $L$, we obtain
\[
\frac1L\left|\sum_{a_-\le n<a_+}e(q(n))\right|
\lesssim_{\mathfrak a_d}
\left(
(\mathcal F L^{-d})^{1/(P_d-2)}
+\mathcal F^{-1}
\right).
\]
%where the fixed factors involving $a_*$ are included in the
%constant.

Passing from $[a_-,a_+)$ to $[a_-,a_+]$ adds at most one summand
of absolute value one. Hence,
\[
\frac1L\left|\sum_{n\in I\cap\Z}e(q(n))\right|
\lesssim_{\mathfrak a_d}
\left(
L^{-1}
+(\mathcal F L^{-d})^{1/(P_d-2)}
+\mathcal F^{-1}
\right).
\]
This proves
\Cref{lem3-add-derivatives}(iii).
%\medskip
%\noindent\textit{Bounded-variation weights.}
%Let $A\le B$ be consecutive integer endpoints in the summation
%subinterval and put $S_n=\sum_{j=A}^ne(q(j))$.
%Finite summation by parts gives
%\[
% \sum_{n=A}^{B}a(n)e(q(n))
% =a(B)S_B+\sum_{n=A}^{B-1}(a(n)-a(n+1))S_n.
%\]
%Every partial sum obeys the corresponding unweighted bound with the
%original scale. Hence the absolute value is at most that bound times
%\[
% |a(B)|+\sum_{n=A}^{B-1}|a(n+1)-a(n)|\le\mathcal V_I(a).
%\]
%This proves the weighted statements without an extra constant.
%Only the specified point values of $a$ enter, so no one-sided
%continuity convention for functions of bounded variation is required.
\end{proof}
\begin{lemma}\label[lemma]{lem3-add-stationary}
Let $I=[u_-,u_+]$, with $u_-<u_+$, and let
$I_a,I_*\subset(u_-,u_+)$ be fixed closed intervals.
Fix positive $A_0,B_0,c_0,C_0$, and denote these data by
$
 \mathfrak s=(I,I_a,I_*,A_0,B_0,c_0,C_0).
$
Suppose $a\in C_c^\infty(\R;\C)$ has support in $I_a$ and
$\|a\|_{C^6(\R)}\le A_0$. Let
$p\in C^\infty(\mathcal O;\R)$ on an open interval
$\mathcal O\supset I$, with
\[
 \|p\|_{C^8(I)}\le B_0,\qquad
 c_0\le|p''(u)|\le C_0\quad(u\in I).
\]
If $p'(u_p)=0$ for a point $u_p\in I_*$, then it
is unique on $I$. For any real $A\ge1$,
\begin{equation}
 \int_I a(u)e(Ap(u))\,du
 =A^{-1/2}e\bigl(Ap(u_p)+\operatorname{sgn}(p'')/8\bigr)
       \frac{a(u_p)}{|p''(u_p)|^{1/2}}+O_{\mathfrak s}(A^{-3/2}).
 \label{eq3-add-stationary-expansion}
\end{equation}
%The error constant is independent of $a,p,u_p,A$ within these
%bounds and independent of the choice of open domain $\mathcal O$.
%The point $u_p$ need not belong to $\operatorname{supp}a$.
\end{lemma}
\begin{proof}
%The sign $\epsilon_p=\operatorname{sgn}(p'')\in\{-1,1\}$ is
%constant on $I$. Thus $p'$ is strictly monotone and has at most one
%zero there. For $u\in I$, define
%\[
% B_p(u)=2\epsilon_p\int_0^1(1-t)p''(u_p+t(u-u_p))\,dt,
% \qquad z_p(u)=(u-u_p)\sqrt{B_p(u)}.
%\]
%Taylor's integral formula gives
%\[
% c_0\le B_p(u)\le C_0,\qquad
% p(u)=p(u_p)+\epsilon_p z_p(u)^2/2.
%\]
%For $0\le s\le6$,
%\[
% B_p^{(s)}(u)=2\epsilon_p\int_0^1(1-t)t^s
%                       p^{(s+2)}(u_p+t(u-u_p))\,dt.
%\]
%The $C^8(I)$-norm bound for $p$ controls these derivatives expressions. By using
%$B_p\ge c_0$ and the above derivatives $B_p^{(s)}$, 
%for $u\ne u_p$,
%\[
% z_p'(u)=\frac{\epsilon_p p'(u)}{z_p(u)},\qquad
% \frac{c_0}{\sqrt{C_0}}\le z_p'(u)\le\frac{C_0}{\sqrt{c_0}}.
%\]
%Indeed, $|p'(u)|$ lies between $c_0|u-u_p|$ and
%$C_0|u-u_p|$, while $|z_p(u)|$ lies between
%$\sqrt{c_0}|u-u_p|$ and $\sqrt{C_0}|u-u_p|$.
%At $u_p$, $z_p'(u_p)=\sqrt{|p''(u_p)|}$, so the bounds extend
%by continuity.
%
%The map $z_p:I\to z_p(I)$ is therefore an increasing bijection.
%For each fixed $p$, its definition extends smoothly to a slightly
%larger open interval, on which $p''$ retains its sign.
%The inverse $v_p:z_p(I)\to I$ is also the restriction of a smooth
%function on an open neighborhood. Only estimates on the displayed
%closed intervals will be used; no common size of the open
%neighborhood is asserted.
%For explicit inverse derivative bounds, put
%\[
% R_{1,p}(u)=1/z_p'(u),\qquad
% R_{s+1,p}(u)=R_{s,p}'(u)/z_p'(u).
%\]
%Induction gives $v_p^{(s)}=R_{s,p}\circ v_p$.
%The uniform lower bound for $z_p'$ and its derivatives through
%order six bound $v_p^{(s)}$ for $1\le s\le5$ by
%$C_{\mathfrak s}$. Thus the transformed amplitude
%\[
% b_{a,p}(z)=a(v_p(z))v_p'(z)\quad(z\in z_p(I)),\qquad
% b_{a,p}(z)=0\quad(z\notin z_p(I))
%\]
%satisfies $\|b_{a,p}\|_{C^4(\R)}\le C_{\mathfrak s}$.
%Its support lies in the fixed interval
%$[-\sqrt{C_0}|I|,\sqrt{C_0}|I|]$.
%The zero extension is smooth: the support of $a$ is separated
%from the endpoints of $I$ by a distance fixed by $I,I_a$, and the
%lower bound for $z_p'$ preserves a positive separation from the
%endpoints of $z_p(I)$.
%Moreover,
%\[
% \begin{gathered}
% b_{a,p}(0)=a(u_p)|p''(u_p)|^{-1/2},\\
% \int_Ia(u)e(Ap(u))\,du
% =e(Ap(u_p))\int_\R b_{a,p}(z)e(\epsilon_pAz^2/2)\,dz.
% \end{gathered}
%\]
%
%Apply \cite[Theorem~2.2, author's version]{Blair2022} in dimension
%one, with expansion order one, scalar matrix $\epsilon_p$, and
%small parameter $\delta=(2\pi A)^{-1}$.
%The normalization $e(t)=\exp(2\pi it)$ gives
%\[
% \begin{split}
% \int_\R b_{a,p}(z)e(\epsilon_pAz^2/2)\,dz
% &=A^{-1/2}e(\epsilon_p/8)b_{a,p}(0)+\mathcal R_A(a,p),\\
% |\mathcal R_A(a,p)|
% &\le C_{\rm abs}A^{-3/2}
%                 \sum_{s=0}^3\|b_{a,p}^{(s)}\|_{L^2(\R)}.
% \end{split}
%\]
%Each $L^2$ norm is bounded by its supremum norm times the square
%root of the fixed support length. Hence
%$|\mathcal R_A(a,p)|\le C_{\mathfrak s}A^{-3/2}$.
%Substituting the value of $b_{a,p}(0)$ proves
%\eqref{eq3-add-stationary-expansion}. This argument is unchanged
%when $a(u_p)=0$, including passage of the critical point across
%the amplitude support.
%\begin{proof}
%	All constants denoted by $C_{\mathfrak s}$ depend only on the
%	fixed data $\mathfrak s$ in \Cref{lem3-add-stationary}.
The proof is divided into three steps.

	\medskip
	\noindent\textbf{Step 1. Reduction to a quadratic phase.}
	Since $p''$ is continuous and never zero on $I$, its sign is
	constant. Put
	$
	\epsilon_p=\operatorname{sgn}(p''(u_p))\in\{-1,1\}.
$
	$p'$ is strictly monotone, so $u_p$ is its unique
	zero on $I$. For $u\in I$, define
	\[
	B_p(u)=2\epsilon_p\int_0^1(1-t)
	p''(u_p+t(u-u_p))\,dt,
	\qquad
	z_p(u)=(u-u_p)\sqrt{B_p(u)}.
	\]
	Taylor's formula and $p'(u_p)=0$ give
	\[
	c_0\le B_p(u)\le C_0,
	\qquad
	p(u)=p(u_p)+\frac{\epsilon_p}{2}z_p(u)^2.
	\]
	%Thus the change of variables is exact.
	
	Differentiation gives
	\[
	z_p'(u)=
	\frac{\epsilon_p\displaystyle\int_0^1
		p''(u_p+t(u-u_p))\,dt}
	{\sqrt{B_p(u)}}.
	\]
	This identity also holds at $u=u_p$, where
	$z_p'(u_p)=\sqrt{|p''(u_p)|}$. Hence,
	\[
	\frac{c_0}{\sqrt{C_0}}
	\le z_p'(u)\le
	\frac{C_0}{\sqrt{c_0}}
	\qquad(u\in I).
	\]
	In particular, $z_p$ is strictly increasing. Write
	\[
	J_p=(z_p(u_-),z_p(u_+)),
	\qquad
	v_p:J_p\longrightarrow(u_-,u_+)
	\]
	for its smooth inverse.
	
	\medskip
	\noindent\textbf{Step 2. Bounds for the transformed amplitude.}
	For $0\le r\le6$,
	\[
	B_p^{(r)}(u)
	=2\epsilon_p\int_0^1(1-t)t^r
	p^{(r+2)}(u_p+t(u-u_p))\,dt.
	\]
	The bound $\|p\|_{C^8(I)}\le B_0$ therefore bounds all these
	derivatives. Since $B_p\ge c_0$, the definition of $z_p$ gives
	\[
	\|z_p\|_{C^6(I)}\lesssim_{\mathfrak s} 1.
	\]
	The inverse identity yields
	\[
	v_p'(z)=\frac1{z_p'(v_p(z))},
	\qquad
	v_p''(z)=-
	\frac{z_p''(v_p(z))}{(z_p'(v_p(z)))^3}.
	\]
	Repeated differentiation of the first identity, using the lower
	bound for $z_p'$ and the bounds for its higher derivatives, gives
	\[
	\max_{1\le r\le5}\sup_{z\in J_p}|v_p^{(r)}(z)|
	\lesssim_{\mathfrak s} 1.
	\]
	
	Define
	\[
	b_{a,p}(z)=
	\begin{cases}
	a(v_p(z))v_p'(z),&z\in J_p,\\
	0,&z\notin J_p.
	\end{cases}
	\]
	Its support is contained in $z_p(I_a)$, a compact subset of
	$J_p$. Thus $b_{a,p}$ is smooth. The product rule,
	the chain rule, and the preceding bounds imply
	\[
	\|b_{a,p}\|_{C^4(\R)}\lesssim_{\mathfrak s} 1.
	\]
	Moreover, $|z_p(u)|\le\sqrt{C_0}|I|$ for $u\in I$, so
	\[
	\operatorname{supp}b_{a,p}
	\subset[-\sqrt{C_0}|I|,\sqrt{C_0}|I|],
	\qquad
	\sum_{r=0}^{3}\|b_{a,p}^{(r)}\|_{L^2(\R)}
	\lesssim_{\mathfrak s} 1.
	\]
	Since $z_p(u_p)=0$, we also have
	$
	b_{a,p}(0)={a(u_p)}/{\sqrt{|p''(u_p)|}}.
	$
	
	\medskip
	\noindent\textbf{Step 3. Evaluation of the quadratic integral.}
	The substitution $z=z_p(u)$ gives
	\[
	\int_I a(u)e(Ap(u))\,du
	=
	e(Ap(u_p))
	\int_\R b_{a,p}(z)e(\epsilon_p A z^2/2)\,dz.
	\]
	Apply \cite[Theorem~2.2]{Blair2022} in
	dimension one, with $k=1$, scalar matrix
	$\epsilon_p$, and $\delta=(2\pi A)^{-1}$. It gives
	\[
	\begin{aligned}
	\int_\R b_{a,p}(z)e(\epsilon_p A z^2/2)\,dz
	&=A^{-1/2}e(\epsilon_p/8)b_{a,p}(0)
	+\mathcal R_A(a,p),\\
	|\mathcal R_A(a,p)|
	&\lesssim A^{-3/2}
	\sum_{r=0}^{3}\|b_{a,p}^{(r)}\|_{L^2(\R)}
	\lesssim_{\mathfrak s}A^{-3/2}.
	\end{aligned}
	\]
	Substituting the value of $b_{a,p}(0)$ proves this lemma.
%\end{proof}
\end{proof}
To conclude this section, we give a remark on nonstationary integration by parts.

\medskip
\noindent\textbf{Remark.}
Fix a closed interval $I=[u_-,u_+]$ with $u_-<u_+$, a closed
interval $K\subset(u_-,u_+)$, an integer $B\ge1$, and positive
constants $c_0,A_B,P_2,\ldots,P_{B+1}$. Write
$$
 \mathfrak n_B=(I,K,B,c_0,A_B,P_2,\ldots,P_{B+1}).
$$
Suppose $a\in C_c^\infty(\R;\C)$ has support in $K$ and
$\|a\|_{C^B(\R)}\le A_B$. Let $p$ be real and smooth on an
open interval $\mathcal O\supset I$, with
$|p'|\ge c_0$ and $|p^{(s)}|\le P_s$ on $I$ for
$2\le s\le B+1$. At this point, $I$ lies in one open component $\mathcal O_p$
of $\{u\in\mathcal O:p'(u)\ne0\}$.
On that component, define
\[
a_{p,0}=a,\qquad a_{p,j+1}
=-\frac d{du}(a_{p,j}/p')\quad(0\le j<B).
\]
Then we have
\begin{equation}
 \int_Ia(u)e(Ap(u))\,du
 =(2\pi iA)^{-B}\int_Ia_{p,B}(u)e(Ap(u))\,du
 =O_{\mathfrak n_B}(A^{-B})\qquad(A\ge1).
 \label{eq3-add-nonstationary}
\end{equation}
%No bound for $p$ or for the upper size of $p'$ is required.
%The constant is also independent of the choice of $\mathcal O$.
%Here is the definition of $a_{p,B}$ and the proof of its uniform
%bound. The interval $I$ lies in one open component $\mathcal O_p$
%of $\{u\in\mathcal O:p'(u)\ne0\}$.
%On that component define
%\[
% a_{p,0}=a,\qquad a_{p,j+1}
%       =-\frac d{du}(a_{p,j}/p')\quad(0\le j<B).
%\]
%These functions retain support in $K$ and extend smoothly by zero
%outside $\mathcal O_p$. Integration by parts introduces no
%endpoint term, since the amplitude vanishes near the endpoints of
%$I$. 

\noindent \textit{Proof of \eqref{eq3-add-nonstationary}.} 
%Iterating gives the exact equality in
%\eqref{eq3-add-nonstationary}.
Since
\[
\frac{d}{du}e(Ap(u))=2\pi iA\,p'(u)e(Ap(u)),
\]
one integration by parts gives
\[
\begin{aligned}
\int_I a(u)e(Ap(u))\,du
&=\frac1{2\pi iA}
\int_{u_-}^{u_+}\frac{a(u)}{p'(u)}
\frac{d}{du}e(Ap(u))\,du\\
&=\frac1{2\pi iA}
\left[\frac{a(u)}{p'(u)}e(Ap(u))\right]_{u_-}^{u_+}
-\frac1{2\pi iA}
\int_I\left(\frac{a(u)}{p'(u)}\right)'e(Ap(u))\,du\\
&=\frac1{2\pi iA}\int_I a_{p,1}(u)e(Ap(u))\,du.
\end{aligned}
\]
%The boundary term vanishes because $a$ is supported in
%$K\Subset(u_-,u_+)$. Here
%\[
%a_{p,0}=a,\qquad
%a_{p,j+1}
%=-\left(\frac{a_{p,j}}{p'}\right)'
%\qquad(0\le j<B).
%\]
%Division by the nonvanishing function $p'$ and differentiation
%preserve support in $K$. Thus every $a_{p,j}$ vanishes near
%both endpoints of $I$. 
Repeating the same integration by parts
$B-1$ more times gives
\[
\int_I a(u)e(Ap(u))\,du
=\frac1{(2\pi iA)^B}
\int_I a_{p,B}(u)e(Ap(u))\,du,
\]
which is the exact equality in \eqref{eq3-add-nonstationary}.

For the estimate, set $Q_p=1/p'$ on $\mathcal O_p$.
Differentiation of $p'Q_p=1$ gives
\[
 Q_p^{(r)}=-\frac1{p'}\sum_{s=1}^r
              \binom rs p^{(s+1)}Q_p^{(r-s)}\qquad(1\le r\le B).
\]
Starting with $|Q_p(u)|\le c_0^{-1},u\in I$, the above recurrence bounds
$\|Q_p\|_{C^B(I)}$ using only $\mathfrak n_B$.
The recurrence for $a_{p,j}$ then gives
$\|a_{p,B}\|_{L^{1}(\R)}\lesssim_{\mathfrak n_B} 1$.
This proves \eqref{eq3-add-nonstationary}.\hfill$\square$

%\clearpage
\section{A smooth Poisson transformation}\label[appendix]{app:oscillatory}

\Cref{lem3-add-dual-formula} isolates the parameter-uniform
Poisson formula used in \Cref{prop3:diagonal}.
%The proof keeps the amplitude error, phase error, and absolute
%summability separate.

\begin{lemma}\label[lemma]{lem3-add-dual-formula}
Let $J=[v_-,v_+]$ with $v_-<v_+$ and let
$J_0\subset(v_-,v_+)$ be a closed interval. Fix
$\Omega\in C_c^\infty((1/2,2);\C)$ and $A_0,B_0,b_0>0$;
write
$
 \mathfrak p=(J,J_0,\Omega,A_0,B_0,b_0).
$
Suppose $a\in C_c^\infty(\R;\C)$ and
$\Psi\in C^\infty(\R;\R)$ satisfy
\[
 \operatorname{supp}a\subset J_0,\quad
 \|a\|_{C^2(\R)}\le A_0,\quad
 \sup_\R|\Psi^{(s)}|\le B_0\ (1\le s\le3),\quad
 \inf_J|\Psi'|\ge b_0.
\]
For real $M\ge1$, $\tau>0$, and $\xi\in\R$, define
\[
 \mathfrak D_{M,\tau}(a,\Psi;\xi)
 =\sum_{m\in\Z}\Omega(m/M)
        \int_\R a(v)e(\xi\Psi(v)-\tau mv)\,dv.
\]
For each fixed $C_*>0$, uniformly for $\tau\ge1$ and
$|\xi|\le C_*M\tau$,
\begin{equation}
 \begin{split}
 \mathfrak D_{M,\tau}(a,\Psi;\xi)
 =\frac1\tau\sum_{r\in\Z}a(r/\tau)
   \Omega\left(\frac{\xi\Psi'(r/\tau)}{M\tau}\right)
                                  e(\xi\Psi(r/\tau))
 +O_{\mathfrak p,C_*}((M\tau)^{-1}).
 \end{split}
 \label{eqB:poisson-formula}
\end{equation}
Also,
\begin{align}
 \sup_{\xi\in\R}|\mathfrak D_{M,\tau}(a,\Psi;\xi)|
 &\lesssim_{\mathfrak p}1+\tau^{-1}
                            &&(M\ge1,\ \tau>0),
       \label{eqB:absolute}\\
 |\mathfrak D_{M,\tau}(a,\Psi;\xi)|
 &\lesssim_{\mathfrak p}(M\tau)^{-1}
             &&(\tau\ge1,\ |\xi|>C_\infty M\tau),
       \label{eqB:large-frequency}
\end{align}
where $C_\infty=\max\{1,4/b_0\}$. 
%The constants depend only on
%$\mathfrak p$, with the additional dependence on $C_*$ in
%\eqref{eqB:poisson-formula}; they are independent of
%$a,\Psi,M,\tau,\xi$ within the stated bounds.
%A nonzero leading summand in \eqref{eqB:poisson-formula} requires
%\[
% r/\tau\in J_0,\qquad
% \frac{M\tau}{2B_0}<|\xi|<\frac{2M\tau}{b_0},\qquad
% \xi\Psi'(r/\tau)>0.
%\]
\end{lemma}
\begin{proof}
The proof is divided into three steps.
	
\medskip
\noindent\textbf{Step 1. Periodization and the absolute bound.}
Extend $\Omega$ by zero to $\R$ and define the $1$-periodic
trigonometric polynomial
$$B_M(z)=\sum_{m\in\Z}\Omega(m/M)e(-mz)$$ for real $z$.
By Poisson summation,
\begin{equation}
 B_M(z)=M\sum_{r\in\Z}\widehat\Omega(M(z+r))=M\sum_{r\in\Z}\widehat\Omega(M(z-r)).
 \label{eq3-add-periodic-poisson}
\end{equation}
%For the sign, the Fourier transform of
%$$s\mapsto\Omega(s/M)e(-sz)$$ at an integer $r$ equals
%$M\widehat\Omega(M(z+r))$; replace $r$ by $-r$ in its integer sum.
For fixed $M$, rapid decay of $\widehat\Omega$ gives absolute and
locally uniform convergence on the right.
Then
\[
 \begin{split}
 \int_0^1|B_M(z)|\,dz
 \le\sum_{r\in\Z}\int_0^1M|\widehat\Omega(M(z-r))|\,dz
 =\sum_{r\in\Z}\int_{-Mr}^{M(1-r)}|\widehat\Omega(s)|\,ds
 =\|\widehat\Omega\|_{L^{1}(\R)}.
 \end{split}
\]
%The last intervals partition $\R$ up to endpoints.
The interval $\tau J_0$ is covered by at most $\tau|J_0|+2$
unit intervals with integer endpoints. Therefore, by periodicity,
\[
 \begin{split}
 |\mathfrak D_{M,\tau}(a,\Psi;\xi)|
 \le\frac{A_0}{\tau}\int_{\tau J_0}|B_M(z)|\,dz
 \le A_0\left(|J_0|+2/\tau\right)\|\widehat\Omega\|_{L^{1}(\R)}.
 \end{split}
\]
This proves \eqref{eqB:absolute}.

%\medskip
%\noindent\textbf{Step 2. An exact identity and absolute convergence.}
%To justify the interchange, first make the same substitution in
%the integral of each absolute value. For a fixed $z\in\R$, a
%nonzero term must satisfy
%$r\in\tau J_0-z/M$, so there are at most $\tau|J_0|+2$
%possibilities. The sum of absolute integrals is consequently at most
%\[
% \frac{A_0(\tau|J_0|+2)}{\tau}
%                   \int_\R|\widehat\Omega(z)|\,dz<\infty.
%\]
%This verifies absolute integrability before as well as after the
%substitution and proves \eqref{eq3-add-dual-exact} by Tonelli and
%Fubini.

\medskip
\noindent\textbf{Step 2. Taylor expansion with an explicit lattice count.}
Fix $C_{*}>0$ and fix $M\ge 1,\tau\ge 1,|\xi|\le C_{*}M\tau$. 
%For the asymptotic formula, assume $\tau\ge1$.
By \eqref{eq3-add-periodic-poisson}
and the substitution $z=M(\tau v-r)$,
\begin{equation}
\mathfrak D_{M,\tau}(a,\Psi;\xi)
=\frac1\tau\sum_{r\in\Z}\int_\R\widehat\Omega(z)
a\left(\frac{r+z/M}{\tau}\right)
e\left(\xi\Psi\left(\frac{r+z/M}{\tau}\right)\right)\,dz.
\label{eq3-add-dual-exact}
\end{equation}
The sum of the absolute integrals, multiplied by $1/\tau$, is at most
$A_0(|J_0|+2/\tau)\|\widehat\Omega\|_{L^{1}(\R)}$, since a nonzero
summand requires $r\in\tau J_0-z/M$; this justifies the interchange.

For each integer
$r$ and real $z$, write
\[
 s_r=r/\tau,\qquad q_z=z/(M\tau).
\]
%Their dependence on the fixed scale parameters is understood; the
%indices indicate which variable changes in the subsequent sum or
%integral. 
Replace the factor
$a(s_r+q_z)e(\xi\Psi(s_r+q_z))$ by
$a(s_r)e(\xi\Psi(s_r)+\xi\Psi'(s_r)q_z)$.
The exact difference is $E_1(r,z)+E_2(r,z)$, where
\begin{align*}
 E_1(r,z)&=(a(s_r+q_z)-a(s_r))e(\xi\Psi(s_r+q_z)),\\
 E_2(r,z)&=a(s_r)\left[e(\xi\Psi(s_r+q_z))
             -e(\xi\Psi(s_r)+\xi\Psi'(s_r)q_z)\right].
\end{align*}
The mean value bound for $a$ and the integral remainder formula
\[
 \Psi(s_r+q_z)-\Psi(s_r)-q_z\Psi'(s_r)
 =q_z^2\int_0^1(1-t)\Psi''(s_r+tq_z)\,dt
\]
give
\[
 |E_1(r,z)|\le A_0|q_z|,\qquad
 |E_2(r,z)|\le 2\pi A_0B_0|\xi|q_z^2.
\]
%These inequalities hold for every real $q_z$; the Taylor formula
%uses the globally defined function $\Psi$ and its global derivative
%bound.

For fixed $z$, $E_1(r,z)$ is zero unless
$s_r\in J_0\cup(J_0-q_z)$, and $E_2(r,z)$ is zero unless
$s_r\in J_0$. Each of these intervals contains at most
$\tau|J_0|+2$ points of $\tau^{-1}\Z$.
%Counting the union, rather than the interval between its extreme
%endpoints, makes the bound independent of the magnitude of $q_z$.
It follows that
\begin{equation}
 \begin{split}
 \sum_{r\in\Z}(|E_1(r,z)|+|E_2(r,z)|)
% \le2A_0(\tau|J_0|+2)\frac{|z|}{M\tau}
% +2\pi A_0B_0(\tau|J_0|+2)
%                         \frac{|\xi|z^2}{M^2\tau^2}
 \lesssim_{\mathfrak p}(\tau+1)
       \left(\frac{|z|}{M\tau}+\frac{|\xi|z^2}{M^2\tau^2}\right).
 \end{split}
 \label{eq3-add-dual-error}
\end{equation}
Let $$\mu_i(\Omega)=\int_\R|z|^i|\widehat\Omega(z)|\,dz$$ for
$i=1,2$. Both are finite since $\Omega$ is smooth with compact
support. Integrating \eqref{eq3-add-dual-error} against
$|\widehat\Omega(z)|/\tau$ gives an error at most
\[
 C_{\mathfrak p}(1+\tau^{-1})
 \left(\frac{\mu_1(\Omega)}{M\tau}
       +\frac{|\xi|\mu_2(\Omega)}{M^2\tau^2}\right)
 \le\frac{C_{\mathfrak p,C_*}}{M\tau}.
\]
%because $\tau\ge1$ and $|\xi|\le C_*M\tau$.

Only finitely many $r$ have $s_r\in J_0$.
Fourier inversion in each of them gives
\[
 \int_\R\widehat\Omega(z)
       e\left(\frac{\xi\Psi'(s_r)}{M\tau}z\right)\,dz
 =\Omega\left(\frac{\xi\Psi'(s_r)}{M\tau}\right).
\]
This proves \eqref{eqB:poisson-formula}.
%and its stated error bound.
%For a nonzero leading summand,
%$s_r\in J_0$ and $\xi\Psi'(s_r)/(M\tau)\in(1/2,2)$.
%The inequalities $b_0\le|\Psi'(s_r)|\le B_0$ give exactly the
%frequency restrictions in \Cref{lem3-add-dual-formula}.

\medskip
\noindent\textbf{Step 3. The complementary large-frequency range.}
Suppose $|\xi|>C_\infty M\tau$ and $\tau\ge1$.
For each relevant integer $m$, $|m|<2M$, and for $v\in J$,
\[
 |\xi\Psi'(v)-\tau m|
 \ge b_0|\xi|-2M\tau\ge b_0|\xi|/2.
\]
Normalize the phase by setting
\[
 p_{m,\tau,\xi}(v)=\operatorname{sgn}(\xi)\Psi(v)
                          -\frac{\tau m}{|\xi|}v,\qquad v\in\R.
\]
Its first derivative has absolute value at least $b_0/2$ on $J$,
and its second and third derivatives there are bounded by $B_0$.
The amplitude has support in $J_0$ and a $C^2$ bound $A_0$.

Applying \eqref{eq3-add-nonstationary} with $B=2$ and large
parameter $|\xi|>1$ therefore bounds each integral $$\int_\R a(v)e(\xi\Psi(v)-\tau mv)\,dv$$ by
$C_{\mathfrak p}|\xi|^{-2}$. There are at most $3M$ relevant
integers $m$, so
\[
 |\mathfrak D_{M,\tau}(a,\Psi;\xi)|
 \lesssim_{\mathfrak p}M|\xi|^{-2}
 \lesssim_{\mathfrak p}{(M\tau)^{-1}}.
\]
This proves \eqref{eqB:large-frequency} and completes the proof of
\Cref{lem3-add-dual-formula}.
\end{proof}

%\clearpage
\section{An arithmetic counting estimate}\label[appendix]{app:resonance}
\Cref{lem3-add-resonance} isolates the counting estimate used in \Cref{prop3:offdiagonal}.
\begin{lemma}\label[lemma]{lem3-add-resonance}
Fix $0<c_1<c_2$ and positive $\tau_*,q_*,B,C_*$, and write
$
 \mathfrak r=(c_1,c_2,\tau_*,q_*,B,C_*).
$
Let $N\ge2$ be an integer and let $Q,A,\tau>0$ satisfy
\[
 Q\ge q_*,\qquad A=\tau Q,\qquad \tau\ge\tau_*,\qquad
 A+Q\le C_*N^B.
\]
Define $\mathcal P_Q=[c_1Q,c_2Q]\cap\N$ and, for $p\in\mathcal P_Q$,
let $d_A(p)=\|A/p\|_{\R/\Z}\in[0,1/2]$.
For every $\epsilon>0$ and $v\in(0,1/2]$,
\begin{align}
 |\{p\in\mathcal P_Q:d_A(p)\le v\}|
 &\lesssim_{\mathfrak r,\epsilon}N^\epsilon(1+Qv),
       \label{eq3-add-resonance-count}\\
 \sum_{\substack{p\in\mathcal P_Q\\d_A(p)>v}}d_A(p)^{-1}
 &\lesssim_{\mathfrak r,\epsilon}
         N^\epsilon\bigl(v^{-1}+Q\log(2/v)\bigr).
       \label{eq3-add-resonance-reciprocal}
\end{align}
%The constants are independent of $N,Q,A,\tau,v$. In particular,
%$A$ is a positive real parameter, not necessarily an integer.
\end{lemma}
\begin{proof}
For $p\in\mathcal P_Q$ with $d_A(p)\le v$, choose the nonnegative
nearest integer
\[
 \ell_p(A)=\lfloor A/p+1/2\rfloor.
\]
%This convention resolves a tie toward the larger integer.
%Since $A/p>0$, no negative integer is needed, and
Then
\[
 |A-p\ell_p(A)|=p\,d_A(p)\le pv\le c_2Qv.
\]
We first count those indices with $\ell_p(A)\ge1$.
$b_p(A):=p\ell_p(A)$ are positive integers in
$[A-c_2Qv,A+c_2Qv]$, an interval containing at most $2c_2Qv+2$
integers. Moreover,
\[
 b_p(A)\le A+c_2Q/2
 \le\max\{1,c_2/2\}(A+Q)
 \le C_1^{\rm size}N^B,
 \qquad C_1^{\rm size}=\max\{1,c_2/2\}C_*.
\]

For each possible product $b$, every eligible $p$ is a positive
divisor of $b$. We may bound their number by the full divisor
function $d_{\rm div}(b)$.
%For clarity, its needed uniform estimate follows directly from
%prime factorization. 
Write $$b=\prod_{\ell\ \textup{prime}}\ell^{v_\ell(b)}$$ and fix $\delta>0$. For prime $\ell$ with $\ell^\delta\ge2$,
\[
 a+1\le2^a\le\ell^{\delta a}\qquad(a\in\Z_{\ge0}).
\]
For each of the finitely many other primes,
$$C_{\ell,\delta}:=\sup_{a\in\Z_{\ge0}}((a+1)\ell^{-\delta a})$$
is finite. Multiplying gives
\[
 d_{\rm div}(b)=\prod_{\ell\ \textup{prime}}(v_\ell(b)+1)
 \le\left(\prod_{\substack{\ell^\delta<2\\ \ell\ \textup{prime}}}C_{\ell,\delta}\right)b^\delta
 =C_\delta b^\delta.
\]
The constant $C_\delta$ depends on $\delta$ alone; the same inequality
includes $b=1$. Choose $\delta=\epsilon/(2B)$. For
$b\le C_1^{\rm size}N^B$,
\[
 d_{\rm div}(b)\le C_\delta(C_1^{\rm size})^\delta N^{\epsilon/2}
                  \lesssim_{\mathfrak r,\epsilon}N^\epsilon.
\]
Multiplication by the number of possible products proves a total
bound $C_{\mathfrak r,\epsilon}N^\epsilon(1+Qv)$ for
$\ell_p(A)\ge1$.

If $\ell_p(A)=0$, then $A/p=d_A(p)\le v$. Hence
\[
 v\ge\frac A p\ge\frac\tau{c_2}\ge\frac{\tau_*}{c_2}.
\]
The total number of such candidates is at most
\[
 |\mathcal P_Q|\le(c_2-c_1)Q+1
 \le(c_2-c_1+q_*^{-1})Q
 \le\frac{c_2(c_2-c_1+q_*^{-1})}{\tau_*}\,Qv.
\]
Adding this bound proves \eqref{eq3-add-resonance-count}.
%No divisor estimate is applied to a zero product.

For \eqref{eq3-add-resonance-reciprocal}, the sum is empty when
$v=1/2$. If $0<v<1/2$, define, for integers $j\ge0$ with
$2^jv<1/2$,
\[
 w_j(v)=\min\{2^{j+1}v,1/2\},\qquad
 E_j(v)=\{p\in\mathcal P_Q:2^jv<d_A(p)\le w_j(v)\}.
\]
These are disjoint sets covering all indices with $d_A(p)>v$,
including $d_A(p)=1/2$. Applying
\eqref{eq3-add-resonance-count} at $w_j(v)\le1/2$ gives
\[
 \begin{split}
 \sum_{p\in E_j(v)}d_A(p)^{-1}
 \le(2^jv)^{-1}C_{\mathfrak r,\epsilon}N^\epsilon
                                    (1+Qw_j(v))
 \le C_{\mathfrak r,\epsilon}N^\epsilon
                                    (2^{-j}v^{-1}+2Q).
 \end{split}
\]
The first terms sum to at most
$2C_{\mathfrak r,\epsilon}N^\epsilon v^{-1}$.
The number of admissible indices $j$ is at most
\[
 1+\log_2(1/(2v))\lesssim \log(2/v).
\]
The second terms therefore sum to at most
$C_{\mathfrak r,\epsilon}N^\epsilon Q\log(2/v)$.
This proves \eqref{eq3-add-resonance-reciprocal} and
\Cref{lem3-add-resonance}.
\end{proof}
%\clearpage
\begin{thebibliography}{99}

\bibitem{BLM2012}
V. Bergelson, A. Leibman, and C. G. Moreira.
\newblock From discrete- to continuous-time ergodic theorems.
\newblock \emph{Ergodic Theory Dynam. Systems} \textbf{32} (2012), no.~2, 383--426.

\bibitem{Blair2022}
M. D. Blair.
\newblock The Van Vleck formula on Ehrenfest time scales and stationary
  phase asymptotics for frequency-dependent phases.
\newblock \emph{Comm. Math. Phys.} \textbf{392} (2022), 517--543.
%\newblock \href{https://doi.org/10.1007/s00220-022-04384-z}{doi:10.1007/s00220-022-04384-z}.
%\newblock Author's version: \href{https://arxiv.org/abs/2012.13034v2}{arXiv:2012.13034v2}.

\bibitem{Boshernitzan1994}
M. D. Boshernitzan.
\newblock Uniform distribution and Hardy fields.
\newblock \emph{J. Anal. Math.} \textbf{62} (1994), 225--240.

\bibitem{Bourgain90}
J. Bourgain.
\newblock Double recurrence and almost sure convergence.
\newblock \emph{J. Reine Angew. Math.} \textbf{404} (1990), 140--161.

\bibitem{Daskalakis2025}
L. Daskalakis.
\newblock Pointwise ergodic theorems for non-conventional bilinear averages
  along {$(\lfloor n^c\rfloor,\allowbreak-\lfloor n^c\rfloor)$}.
\newblock {arXiv:\allowbreak2503.03976}, 2025.

\bibitem{Folland2016}
G. B. Folland.
\newblock \emph{A Course in Abstract Harmonic Analysis}, second edition.
\newblock Textbooks in Mathematics, CRC Press, Boca Raton, FL, 2016.

\bibitem{FornalKrause2026}
J. Fornal and B. Krause.
\newblock Double recurrence and almost sure convergence: Primes and weighted
  theory.
\newblock {arXiv:\allowbreak2603.22203v1}, 2026.

\bibitem{FN2023}
N. Frantzikinakis.
\newblock Joint ergodicity of sequences.
\newblock \emph{Adv. Math.} \textbf{417} (2023), article 108918.

\bibitem{HeathBrown2017}
D. R. Heath-Brown.
\newblock A new $k$th derivative estimate for exponential sums via
  Vinogradov's mean value.
\newblock \emph{Proc. Steklov Inst. Math.} \textbf{296} (2017), 88--103.
%\newblock \href{https://doi.org/10.1134/S0081543817010072}{doi:10.1134/S0081543817010072}.
%\newblock Author's version: \href{https://arxiv.org/abs/1601.04493v3}{arXiv:1601.04493v3}.

\bibitem{HSX2023}
W. Huang, S. Shao, and R. Xiao.
\newblock Pointwise convergence of some continuous-time polynomial ergodic
  averages.
\newblock {arXiv:2310.16780v2}, 2025.

\bibitem{Krause2025}
B. Krause.
\newblock A unified approach to two pointwise ergodic theorems: Double
  recurrence and return times.
\newblock {arXiv:\allowbreak2501.06877}, 2025.

\bibitem{LST2025}
C. Lutsko, A. Sourmelidis, and N. Technau.
\newblock Pair correlation of the fractional parts of $\alpha n^\theta$.
\newblock \emph{J. Eur. Math. Soc.} \textbf{27} (2025), 4069--4082.
%\newblock \href{https://doi.org/10.4171/JEMS/1449}{doi:10.4171/JEMS/1449}.

\end{thebibliography}
\end{document}



